
% BEGIN S2
 
\par\medskip\noindent\textbf{2 Preliminaries}\par
 
% BEGIN S2.p1
 
In this section, we introduce necessary notation and preliminaries. Further, we recall and adapt the definition of interface automata and related notions from [14].\par
 
% END S2.p1
 
% BEGIN S2.SS0.SSS0.Px1
 
\par\medskip\noindent\textbf{Languages and Automata}\par
 
% BEGIN S2.SS0.SSS0.Px1.p1
 
For an alphabet $\Sigma$, the set $\Sigma^{\infty}:=\Sigma^{*}\cup\Sigma^{\omega}$ contains all finite and infinite words over $\Sigma$. Given a sequence $\sigma\in(\Sigma\cup\Sigma^{\prime})^{\infty}$, we denote with $\sigma|_{\Sigma^{\prime}}$ the projection of $\sigma$ to $\Sigma^{\prime}$, that is, $\sigma|_{\Sigma^{\prime}}\in(\Sigma^{\prime})^{\infty}$ is the sequence obtained from $\sigma$ by removing all letters in $\Sigma\setminus\Sigma^{\prime}$. For a sequence $\sigma\in\Sigma^{\infty}$, we denote with $\sigma[i]\in\Sigma$ the letter of $\sigma$ at the $i$-th position.\par
 
% END S2.SS0.SSS0.Px1.p1
 
% BEGIN S2.SS0.SSS0.Px1.p2
 
A non-deterministic finite automaton (NFA) over an alphabet $\Sigma$ is a tuple $\mathcal{N}=\langle Q,\Sigma,\delta,Q_{0},F\rangle$ with finite set of states $Q$, initial states $Q_{0}\subseteq Q$, transition relation $\delta\subseteq Q\times\Sigma\times Q$, and a set of accepting states $F\subseteq Q$. A run of $\mathcal{N}$ on a finite word $\sigma=a_{1}\ldots a_{n}\in\Sigma^{*}$ is a finite sequence $\rho=\rho_{0}\rho_{1}\ldots\rho_{n}\in Q^{*}$ such that $\rho_{0}\in Q_{0}$ and for every $i<n$, it holds that $\rho_{i+1}\in\delta(\rho_{i},\sigma_{i+1})$. A run $\rho$ is accepting if and only if $\rho_{n}\in F$. A non-deterministic Büchi automaton (NBA) $\mathcal{B}=\langle Q,\Sigma,\delta,Q_{0},B\rangle$ on infinite words instead has a set of Büchi-accepting states $B\subseteq Q$ of which some must be visited infinitely often. A run $\rho$ of $\mathcal{B}$ on $\sigma\in\Sigma^{\omega}$ is defined analogously, and is accepting if and only if for every $i\in\mathbb{N}$ there exists $j\geq i$ such that $\rho[j]\in B$. The language of an automaton $\mathcal{L}(\mathcal{A})$ is the set of words $\sigma$ on which $\mathcal{A}$ has an accepting run. For a universal co-Büchi automaton (UCW) $\mathcal{C}=\langle Q,\Sigma,\delta,Q_{0},B\rangle$, a run $\rho$ of $\mathcal{C}$ on a word $\sigma\in\Sigma^{\omega}$ is accepting if and only if there exists $i\in\mathbb{N}$ such that for all $j\geq i$, $\rho[j]\not\in B$, and $\sigma\in\Sigma^{\omega}$ is only accepted by $\mathcal{C}$ if every run is accepting. We define the size of an automaton $\mathcal{A}$ as the total number of states and transitions, i.e. $|\mathcal{A}|=|Q|+|\delta|$.\par
 
% END S2.SS0.SSS0.Px1.p2
 
% BEGIN S2.SS0.SSS0.Px1.p3
 
We use $\mathbb{R}$ to denote the set of real numbers, and define $\infty,-\infty$ as the elements where $\infty>n>-\infty$ for all $n\in\mathbb{R}$. The set $\mathbb{R}_{-\infty}=\mathbb{R}\cup\{-\infty\}$ then contains the real numbers and $-\infty$. We define conventionally the infimum over the empty set as $\inf(\emptyset)=\infty$.\par
 
% END S2.SS0.SSS0.Px1.p3
 
% END S2.SS0.SSS0.Px1
 
% BEGIN S2.SS0.SSS0.Px2
 
\par\medskip\noindent\textbf{Interface Automata}\par
 
% BEGIN S2.SS0.SSS0.Px2.p1
 
Interface automata [14] are a modeling formalism for specifying the interactions between components and their environment. Components communicate via synchronization on input and output actions. In contrast to the interface automata in [14], we distinguish external inputs $\Sigma^{E}$, which are external to the overall system and broadcast to all components, from inputs $\Sigma^{I}$ that come from system components outside of $A$ and are subject to assumptions specified by $A$. The external inputs enable synchronization between multiple components on the same input action, and are used to model behavior uncontrollable by the system.\par
 
% END S2.SS0.SSS0.Px2.p1
 
% BEGIN Thmtheorem2
 
\par\medskip\noindent\textbf{Definition 2 (Interface Automaton (adapted from [14])).}\par
 
% BEGIN Thmtheorem2.p1
 
An interface automaton $A=\langle V,V^{init},\Sigma^{I},\Sigma^{O},\Sigma^{E},\Sigma^{H},\mathcal{T}\rangle$ is a tuple where:\par
  
\par\noindent  $\blacksquare$ 
% BEGIN S2.I1.i1.p1
 
$V$ is a finite set of states and $V^{init}\subseteq V$ is a set of initial states. We require that $V^{init}$ contains at most one state. If $V^{init}=\emptyset$, then $A$ is called empty.\par
 
% END S2.I1.i1.p1
 
 
\par\noindent  $\blacksquare$ 
% BEGIN S2.I1.i2.p1
 
$\Sigma^{I},\Sigma^{O},\Sigma^{E},\Sigma^{H}$ are mutually disjoint finite sets of input, output, external input and internal actions respectively. Let $\Sigma_{A}:=\Sigma^{I}\cup\Sigma^{O}\cup\Sigma^{E}\cup\Sigma^{H}$ be the set of all actions of $A$.\par
 
% END S2.I1.i2.p1
 
 
\par\noindent  $\blacksquare$ 
% BEGIN S2.I1.i3.p1
 
$\mathcal{T}\subseteq V\times\Sigma_{A}\times V$ is a set of transitions such that for every state $v\in V$ it holds that:\par
  
\par\noindent  – 
% BEGIN S2.I1.i3.I1.i1.p1
 
for every external input action $a\in\Sigma^{E}$, there exists $v^{\prime}\in V$ such that $(v,a,v^{\prime})\in\mathcal{T}$;\par
 
% END S2.I1.i3.I1.i1.p1
 
 
\par\noindent  – 
% BEGIN S2.I1.i3.I1.i2.p1
 
for all $a\in\Sigma^{E}\cup\Sigma^{I}$ and $v^{\prime},v^{\prime\prime}\in V$ with $(v,a,v^{\prime})\in\mathcal{T}$ and $(v,a,v^{\prime\prime})\in\mathcal{T}$ we have $v^{\prime}=v^{\prime\prime}$.\par
 
% END S2.I1.i3.I1.i2.p1
 
  
% END S2.I1.i3.p1
 
  
% END Thmtheorem2.p1
 
% END Thmtheorem2
 
% BEGIN S2.SS0.SSS0.Px2.p2
 
Note that we require $A$ to be input enabled on external input actions $\Sigma^{E}$, i.e., it accommodates any external input in every state. We also require that $A$ is input-deterministic on $\Sigma^{E}\cup\Sigma^{I}$, in order to ensure compositionality of parallel composition [15].\par
 
% END S2.SS0.SSS0.Px2.p2
 
% BEGIN S2.SS0.SSS0.Px2.p3
 
We define the size of $A$ as $|A|=|V|+|\mathcal{T}|$. For $u\in V$, we define the interface automaton $A^{u}=\langle V,\{u\},\Sigma^{I},\Sigma^{O},\Sigma^{E},\Sigma^{H},\mathcal{T}\rangle$ obtained from $A$ by replacing the set of initial states by $\{u\}$.\par
 
% END S2.SS0.SSS0.Px2.p3
 
% BEGIN S2.SS0.SSS0.Px2.p4
 
Let $A=\langle V,V^{init},\Sigma^{I},\Sigma^{O},\Sigma^{E},\Sigma^{H},\mathcal{T}\rangle$ be an interface automaton. An execution of $A$ is an alternating sequence $v_{0},a_{0},v_{1},a_{1},\dots$ of states $v_{i}\in V$ and actions $a_{i}\in\Sigma_{A}$ such that $(v_{i},a_{i},v_{i+1})\in\mathcal{T}$ for all $i\geq 0$. An execution fragment is a finite prefix $v_{0},a_{0},v_{1},a_{1},\dots,v_{n}$ of an execution ending in a state. A state $v\in V$ is reachable in $A$ if there exists an execution fragment starting in some $v_{0}\in V^{init}$ that ends in $v$. We denote with $\mathsf{Reach}(A)\subseteq V$ the set of states reachable in $A$. For $\Sigma^{\prime}\subseteq\Sigma_{A}$, let $\mathsf{Enabled}_{A}(v,\Sigma^{\prime}):=\{a\in\Sigma^{\prime}\mid\exists v^{\prime}.\nobreak\ (v,a,v^{\prime})\in\mathcal{T}\}$ be the set of actions from $\Sigma^{\prime}$ enabled in state $v$.\par
 
% END S2.SS0.SSS0.Px2.p4
 
% BEGIN S2.SS0.SSS0.Px2.p5
 
Two key notions in the theory of interface automata are the composability and product of automata. Whether two interface automata are composable depends on their actions. They must not share input, output and internal actions, but may have joint external inputs. For two composable interface automata, we can construct their product. This is formalized in the definitions below. For the remainder of this section, let $A_{P}=\langle V_{P},V_{P}^{init},\Sigma_{P}^{I},\Sigma_{P}^{O},\Sigma_{P}^{E},\Sigma_{P}^{H},\mathcal{T}_{P}\rangle$ and $A_{Q}=\langle V_{Q},V_{Q}^{init},\Sigma_{Q}^{I},\Sigma_{Q}^{O},\Sigma_{Q}^{E},\Sigma_{Q}^{H},\mathcal{T}_{Q}\rangle$ be interface automata.\par
 
% END S2.SS0.SSS0.Px2.p5
 
% BEGIN Thmtheorem3
 
\par\medskip\noindent\textbf{Definition 3 (Composability [14]).}\par
 
% BEGIN Thmtheorem3.p1
 
Two interface automata $A_{P}$ and $A_{Q}$ are composable if $\Sigma_{P}^{I}\cap\Sigma_{Q}^{I}=\emptyset$, $\Sigma_{P}^{O}\cap\Sigma_{Q}^{O}=\emptyset$, $\Sigma_{P}^{H}\cap\Sigma_{Q}=\emptyset$, and $\Sigma_{Q}^{H}\cap\Sigma_{P}=\emptyset$.\par
 
% END Thmtheorem3.p1
 
% BEGIN Thmtheorem3.p2
 
We define $\mathsf{Shared}(A_{P},A_{Q}):=(\Sigma_{P}^{I}\cap\Sigma_{Q}^{O})\cup(\Sigma_{Q}^{I}\cap\Sigma_{P}^{O})$, and denote with $\mathsf{JointE}(A_{P},A_{Q}):=\Sigma_{P}^{E}\cap\Sigma_{Q}^{E}$ the set of joint external inputs of $A_{P}$ and $A_{Q}$.\par
 
% END Thmtheorem3.p2
 
% END Thmtheorem3
 
% BEGIN Thmtheorem4
 
\par\medskip\noindent\textbf{Definition 4 (Product).}\par
 
% BEGIN Thmtheorem4.p1
 
Let $A_{P}$ and $A_{Q}$ be composable. Their product $A_{P}\otimes A_{Q}=\langle V_{P\otimes Q},V_{P\otimes Q}^{init},\Sigma_{P\otimes Q}^{I},\Sigma_{P\otimes Q}^{O},\Sigma_{P\otimes Q}^{E},\Sigma_{P\otimes Q}^{H},\mathcal{T}_{P\otimes Q}\rangle$ is the interface automaton with the components defined as follows:\par
  
\par\noindent  $\blacksquare$ 
% BEGIN S2.I2.i1.p1
 
$V_{P\otimes Q}:=V_{P}\times V_{Q}$, $V_{P\otimes Q}^{init}=V_{P}^{init}\times V_{Q}^{init}$,\par
 
% END S2.I2.i1.p1
 
 
\par\noindent  $\blacksquare$ 
% BEGIN S2.I2.i2.p1
 
$\Sigma_{P\otimes Q}^{I}=(\Sigma_{P}^{I}\cup\Sigma_{Q}^{I})\setminus\mathsf{Shared}(A_{P},A_{Q}),\ \Sigma_{P\otimes Q}^{O}=(\Sigma_{P}^{O}\cup\Sigma_{Q}^{O})\setminus\mathsf{Shared}(A_{P},A_{Q}),\\
\Sigma_{P\otimes Q}^{E}=\Sigma_{P}^{E}\cup\Sigma_{Q}^{E},\ \Sigma_{P\otimes Q}^{H}=(\Sigma_{P}^{H}\cup\Sigma_{Q}^{H})\cup\mathsf{Shared}(A_{P},A_{Q})$,\par
 
% END S2.I2.i2.p1
 
 
\par\noindent  $\blacksquare$ 
% BEGIN S2.I2.i3.p1
 
$\begin{array}[]{lll}\mathcal{T}_{P\otimes Q}&=\{((v_{P},v_{Q}),\alpha,(v_{P}^{\prime},v_{Q}^{\prime}))&\mid\alpha\in(\mathsf{Shared}(A_{P},A_{Q})\cup\mathsf{JointE}(A_{P},A_{Q}))\\
&&\;\;\;\land(v_{P},\alpha,v_{P}^{\prime})\in\mathcal{T}_{P}\land(v_{Q},\alpha,v_{Q}^{\prime})\in\mathcal{T}_{Q}\}\\
&\cup\{((v_{P},v_{Q}),\alpha,(v_{P}^{\prime},v_{Q}))&\mid\alpha\not\in(\mathsf{Shared}(A_{P},A_{Q})\cup\mathsf{JointE}(A_{P},A_{Q}))\\
&&\;\;\;\land(v_{P},\alpha,v_{P}^{\prime})\in\mathcal{T}_{P}\}\\
&\cup\{((v_{P},v_{Q}),\alpha,(v_{P},v_{Q}^{\prime}))&\mid\alpha\not\in(\mathsf{Shared}(A_{P},A_{Q})\cup\mathsf{JointE}(A_{P},A_{Q}))\\
&&\;\;\;\land(v_{Q},\alpha,v_{Q}^{\prime})\in\mathcal{T}_{Q}\}.\end{array}$\par
 
% END S2.I2.i3.p1
 
  
% END Thmtheorem4.p1
 
% END Thmtheorem4
 
% BEGIN S2.SS0.SSS0.Px2.p6
 
Note that the actions $\mathsf{Shared}(A_{P},A_{Q})$ become internal actions of $A_{P}\otimes A_{Q}$, and $P$ and $Q$ must synchronize on shared and joint external input actions.\par
 
% END S2.SS0.SSS0.Px2.p6
 
% BEGIN S2.SS0.SSS0.Px2.p7
 
Given two composable interface automata $A_{P}$ and $A_{Q}$, a product state $(v_{P},v_{Q})\in V_{P}\times V_{Q}$ is called an illegal state of the product automaton $A_{P\otimes Q}$ if and only if there exists $a\in\mathsf{Shared}(A_{P},A_{Q})$ such that $a\in(\mathsf{Enabled}_{P}(v_{P},\Sigma_{P}^{O})\setminus\mathsf{Enabled}_{Q}(v_{Q},\Sigma_{Q}^{I}))\cup(\mathsf{Enabled}_{Q}(v_{Q},\Sigma_{Q}^{O})\setminus\mathsf{Enabled}_{P}(v_{P},\Sigma_{P}^{I}))$. That is, illegal states are ones in which a shared action can be produced as an output of one of the interface automata but is not allowed as an input by the other one. Let $\mathsf{Illegal}(A_{P},A_{Q})$ be the set of illegal states of $A_{P}\otimes A_{Q}$.\par
 
% END S2.SS0.SSS0.Px2.p7
 
% BEGIN S2.SS0.SSS0.Px2.p8
 
We say that two interface automata $A_{P}$ and $A_{Q}$ are compatible if there exists a way for the rest of the system, supplying the inputs $\Sigma_{P\otimes Q}^{I}$ to ensure that illegal states are avoided. This gives rise to the notion of legal environment for $A_{P}$ and $A_{Q}$, which we recall next.\par
 
% END S2.SS0.SSS0.Px2.p8
 
% BEGIN Thmtheorem5
 
\par\medskip\noindent\textbf{Definition 5 (Legal Environment and Compatibility[14]).}\par
 
% BEGIN Thmtheorem5.p1
 
A legal environment for $(A_{P},A_{Q})$ is a non-empty interface automaton $A_{R}=\langle V_{R},V_{R}^{init},\Sigma_{R}^{I},\Sigma_{R}^{O},\Sigma_{R}^{E},\Sigma_{R}^{H},\mathcal{T}_{R}\rangle$ where:\par
 
% END Thmtheorem5.p1
 
% BEGIN Thmtheorem5.p2
  
\par\noindent  1. 
% BEGIN S2.I3.i1.p1
 
$\Sigma_{R}^{E}=\Sigma_{P\otimes Q}^{E}$, $\Sigma_{R}^{I}=\Sigma_{P\otimes Q}^{O}$, $\Sigma_{R}^{O}=\Sigma_{P\otimes Q}^{I}$, $\Sigma_{R}^{H}=\emptyset$.\par
 
% END S2.I3.i1.p1
 
 
\par\noindent  2. 
% BEGIN S2.I3.i2.p1
 
$A_{R}$ is composable with $A_{P}\otimes A_{Q}$, and $\mathsf{Illegal}(A_{P}\otimes A_{Q},A_{R})=\emptyset$.\par
 
% END S2.I3.i2.p1
 
 
\par\noindent  3. 
% BEGIN S2.I3.i3.p1
 
$\mathsf{Reach}((A_{P}\otimes A_{Q})\otimes A_{R})\cap(\mathsf{Illegal}(A_{P},A_{Q})\times V_{R})=\emptyset$.\par
 
% END S2.I3.i3.p1
 
  
% END Thmtheorem5.p2
 
% BEGIN Thmtheorem5.p3
 
Two interface automata $A_{P}$ and $A_{Q}$ are considered compatible if they are non-empty, composable and there exists a legal environment for $(A_{P},A_{Q})$.\par
 
% END Thmtheorem5.p3
 
% END Thmtheorem5
 
% BEGIN S2.SS0.SSS0.Px2.p9
 
Intuitively, a legal environment $A_{R}$ for $(A_{P},A_{Q})$ represents the remaining system beyond $A_{P}$ and $A_{Q}$. Condition 3 requires that the inputs provided by $A_{R}$ to the two interfaces steer them away from illegal states in the product.\par
 
% END S2.SS0.SSS0.Px2.p9
 
% BEGIN Thmtheorem6
 
\par\medskip\noindent\textbf{Definition 6 (Composition of Interface Automata[14]).}\par
 
% BEGIN Thmtheorem6.p1
 
Given composable interface automata $A_{P}$ and $A_{Q}$, a product state $(v_{P},v_{Q})\in V_{P}\times V_{Q}$ is called compatible if there exists a legal environment for $(A_{P}^{v_{P}},A_{Q}^{v_{Q}})$. Let $\mathsf{Comp}(A_{P},A_{Q})$ be the set of compatible product states. The composition $A_{P}\parallel A_{Q}$ of $A_{P}$ and $A_{Q}$ is defined by restricting $A_{P}\otimes A_{Q}$ to $\mathsf{Comp}(A_{P},A_{Q})$. Formally, $A_{P}\parallel A_{Q}=\langle V_{P\otimes Q}\cap\mathsf{Comp}(A_{P},A_{Q}),V_{P\otimes Q}^{init}\cap\mathsf{Comp}(A_{P},A_{Q}),\Sigma_{P\otimes Q}^{I},\Sigma_{P\otimes Q}^{O},\Sigma_{P\otimes Q}^{E},\\
\Sigma_{P\otimes Q}^{H},\mathcal{T}_{P\otimes Q}\cap(\mathsf{Comp}(A_{P},A_{Q})\times\Sigma_{P\otimes Q}\times\mathsf{Comp}(A_{P},A_{Q}))\rangle$.\par
 
% END Thmtheorem6.p1
 
% END Thmtheorem6
 
% BEGIN S2.SS0.SSS0.Px2.p10
 
Intuitively, the compatible product states are those from which an environment can prevent reaching illegal states. Thus, $A_{P}$ and $A_{Q}$ are compatible if and only if $V_{P\otimes Q}^{init}\subseteq\mathsf{Comp}(A_{P},A_{Q})$. Furthermore, for every $(v_{P},v_{Q})\in\mathsf{Comp}(A_{P},A_{Q})$, all the external input actions $\Sigma_{P\otimes Q}^{E}$ lead to product sates in $\mathsf{Comp}(A_{P},A_{Q})$. If this was not the case, $(v_{P},v_{Q})$ would not be in $\mathsf{Comp}(A_{P},A_{Q})$, since an external input action leading to an illegal state cannot be prevented by a legal environment. If an action $a\in\Sigma_{P\otimes Q}^{O}$ leads from $(v_{P},v_{Q})$ to $\mathsf{Illegal}(A_{P},A_{Q})$, we prune all $a$-transitions from $(v_{P},v_{Q})$. Therefore, for compatible interface automata, the composition $A_{P}\parallel A_{Q}$ is a non-empty interface automaton conforming to Definition 2.\par
 
% END S2.SS0.SSS0.Px2.p10
 
% BEGIN S2.SS0.SSS0.Px2.p11
 
Another concept, crucial for independent implementability of interfaces, is refinement. Refinement of interface automata [14] is defined via alternating simulation, recalled below.\par
 
% END S2.SS0.SSS0.Px2.p11
 
% BEGIN S2.SS0.SSS0.Px2.p12
 
For an interface automaton $A$ and state $v$, let $\varepsilon\text{-}\mathsf{closure}_{A}(v)$ be the set of states of $A$ that can be reached from $v$ using only internal actions from $\Sigma^{H}$. We define:\par
  
\par  \quad  \[\begin{array}[]{lll}\mathsf{ExtEn}^{I,E}_{A}(v)&:=&\{a\mid\forall u\in\varepsilon\text{-}\mathsf{closure}_{A}(v).\nobreak\ a\in\mathsf{Enabled}_{A}(u,\Sigma_{A}^{I}\cup\Sigma_{A}^{E})\}\text{ and }\\
\mathsf{ExtEn}^{O}_{A}(v)&:=&\{a\mid\exists u\in\varepsilon\text{-}\mathsf{closure}_{A}(v).\nobreak\ a\in\mathsf{Enabled}_{A}(u,\Sigma_{A}^{O})\}.\end{array}\]\quad  \quad  
  
to be the sets of externally enabled input and output actions at $v$. Further, for actions $a\in\mathsf{ExtEn}^{I,E}_{A}(v)\cup\mathsf{ExtEn}^{O}_{A}(v)$, let $\mathsf{ExtDest}_{A}(v,a)=\{u^{\prime}\mid\exists(u,a,u^{\prime})\in\mathcal{T}_{A}.\nobreak\ u\in\varepsilon\text{-}\mathsf{closure}_{A}(v)\}$.\par
 
% END S2.SS0.SSS0.Px2.p12
 
% BEGIN Thmtheorem7
 
\par\medskip\noindent\textbf{Definition 7 (Alternating Simulation [14]).}\par
 
% BEGIN Thmtheorem7.p1
 
A binary relation $\succeq\subseteq V_{P}\times V_{Q}$ is an alternating simulation from the interface automaton $A_{Q}$ to the interface automaton $A_{P}$ if for all $v_{P}\in V_{P}$ and $v_{Q}\in V_{Q}$, $v_{P}\succeq v_{Q}$ implies that:\par
  
\par\noindent  1. 
% BEGIN S2.I4.i1.p1
 
$\mathsf{ExtEn}^{I,E}_{P}(v_{P})\subseteq\mathsf{ExtEn}^{I,E}_{Q}(v_{Q})$ and $\mathsf{ExtEn}^{O}_{Q}(v_{Q})\subseteq\mathsf{ExtEn}^{O}_{P}(v_{P})$.\par
 
% END S2.I4.i1.p1
 
 
\par\noindent  2. 
% BEGIN S2.I4.i2.p1
 
For all $a\in\mathsf{ExtEn}^{I,E}_{P}(v_{P})\cup\mathsf{ExtEn}^{O}_{Q}(v_{Q})$ and $v_{Q}^{\prime}\in\mathsf{ExtDest}_{Q}(v_{Q},a)$, there exists a state $v_{P}^{\prime}\in\mathsf{ExtDest}_{P}(v_{P},a)$ such that $v_{P}^{\prime}\succeq v_{Q}^{\prime}$.\par
 
% END S2.I4.i2.p1
 
  
% END Thmtheorem7.p1
 
% END Thmtheorem7
 
% BEGIN S2.SS0.SSS0.Px2.p13
 
The existence of an alternating simulation from $A_{Q}$ to $A_{P}$ guarantees that the interactions of $A_{P}$ are preserved in $A_{Q}$. In particular, $A_{Q}$ does not impose more assumptions on the environment than $A_{P}$, and satisfies all the output restrictions of $A_{P}$. That is, any environment compatible with $A_{P}$ also is compatible with $A_{Q}$.\par
 
% END S2.SS0.SSS0.Px2.p13
 
% BEGIN Thmtheorem8
 
\par\medskip\noindent\textbf{Definition 8 (Interface Automata Refinement [14]).}\par
 
% BEGIN Thmtheorem8.p1
 
For interface automata $A_{P}$ and $A_{Q}$, we say that $A_{Q}$ refines $A_{P}$, written $A_{Q}\preceq A_{P}$, if and only if $\Sigma_{Q}^{E}=\Sigma_{P}^{E}$, $\Sigma_{Q}^{I}\supseteq\Sigma_{P}^{I}$, $\Sigma_{Q}^{O}\subseteq\Sigma_{P}^{O}$, and there exists an alternating simulation $\succeq$ from $A_{Q}$ to $A_{P}$, and states $v_{P}\in V_{P}^{init}$ and $v_{Q}\in V_{Q}^{init}$, such that $v_{P}\succeq v_{Q}$.\par
 
% END Thmtheorem8.p1
 
% END Thmtheorem8
 
% BEGIN S2.SS0.SSS0.Px2.p14
 
Thus, the relation $A_{Q}\preceq A_{P}$ guarantees that $A_{Q}$ must accept at least all inputs $\Sigma_{P}^{E}\cup\Sigma_{P}^{I}$, and may not add any new outputs w.r.t. $\Sigma_{P}^{O}$. The alternating simulation ensures that for any input, the outward behavior of $Q$ matches that of $P$.\par
 
% END S2.SS0.SSS0.Px2.p14
  
% END S2.SS0.SSS0.Px2
 
% END S2


% BEGIN S3
 
\par\medskip\noindent\textbf{3 Reward Interfaces}\par
 
% BEGIN S3.p1
 
We now introduce reward interfaces, the central notion of the framework we propose. They build on the classical interface automata, extending them with an additional quantitative requirement, defined as a reward function that assigns numerical values to sequences of observable actions of the interface. Essentially, a reward function defines a quantitative language [8] over the alphabet of non-internal actions of the given interface automaton.\par
 
% END S3.p1
 
% BEGIN Thmtheorem9
 
\par\medskip\noindent\textbf{Definition 9 (Reward Interface).}\par
 
% BEGIN Thmtheorem9.p1
 
A reward interface is a pair $P=(A_{P},\mathcal{F}_{P})$ where $A_{P}=\langle V_{P},V_{P}^{init},\Sigma_{P}^{I},\Sigma_{P}^{O},\Sigma_{P}^{E},\Sigma_{P}^{H},\mathcal{T}_{P}\rangle$ is an interface automaton and $\mathcal{F}_{P}:(\Sigma^{Obs}_{P})^{\infty}\to\mathbb{R}_{-\infty}$ is a partial function that assigns a value from $\mathbb{R}_{-\infty}$ to sequences over the alphabet $\Sigma^{Obs}_{P}:=\Sigma^{I}_{P}\cup\Sigma^{O}_{P}\cup\Sigma^{E}_{P}$ of non-internal actions of $A_{P}$.\par
 
% END Thmtheorem9.p1
 
% END Thmtheorem9
 
% BEGIN S3.p2
 
Intuitively, $\mathcal{F}_{P}$ expresses a quantitative specification, by associating reward values with sequences in $(\Sigma^{Obs}_{P})^{\infty}$. The reward value of a sequence describes how well the respective observed behavior satisfies the quantitative requirement. Since $\mathcal{F}_{P}$ is part of the interface specification, it is defined in terms of the actions $\Sigma^{Obs}_{P}$ that are visible to the component’s environment. We deliberately do not impose further restrictions on the functions $\mathcal{F}_{P}$, in order to retain full generality of the proposed framework. In Section 6 we discuss possible instantiations, in which the reward functions have a natural finite representation.\par
 
% END S3.p2
 
% BEGIN Thmtheorem10
 
\par\medskip\noindent\textbf{Example 10.}\par
 
% BEGIN Thmtheorem10.p1
 
Continuing from Example 1, we show a possible reward interface $P=(A_{P},\mathcal{F}_{P})$ for component $P$. The restrictions from $A_{P}$ (Figure 1(a)) ensure that output $\mathsf{o_{1}}$ can only occur after input $\mathsf{q}$. The requirement that $\mathsf{o_{1}}$ must occur after $\mathsf{e_{1}}$, and not before, is modeled via the reward function $\mathcal{F}_{P}$, mapping sequences over $\Sigma_{P}^{Obs}=\{\mathsf{e_{1}},\mathsf{e_{2}},\mathsf{p},\mathsf{q},\mathsf{o_{1}}\}$ to values. The reward function $\mathcal{F}_{P}$ follows $\mathcal{F}_{\mathsf{S}}$: We award value $1$ if only one type of input $\mathsf{e_{1}}$ or $\mathsf{e_{2}}$ occurs, and require respectively $\mathsf{o_{1}}$ and $\mathsf{p}$, instead of $\mathsf{o_{2}}$, as $P$ has no knowledge nor control over $\mathsf{o_{2}}$. Similarly, if both $\mathsf{e_{1}}$ and $\mathsf{e_{2}}$ occur, the order of $\mathsf{o_{1}}$ and $\mathsf{p}$ must reflect that to achieve value $\frac{1}{2}$, otherwise the value is lowered to $\frac{1}{4}$. In both cases when $\mathsf{e_{2}}$ occurs, $\mathsf{o_{1}}$ and $\mathsf{p}$ should not be performed infinitely, otherwise the assigned value is $0$.\par
 
% END Thmtheorem10.p1
 
% BEGIN Thmtheorem10.p2
 
Note that while we could for example encode the safety requirement that there is no $\mathsf{o_{1}}$ before the first occurrence of $\mathsf{e_{1}}$ as part of the interface automaton, it would make the automaton structure more complicated. Furthermore, as the internal order of $\mathsf{o_{1}}$ and $\mathsf{e_{1}}$ is less relevant for the interaction with component $Q$, which is not concerned with $\mathsf{o_{1}}$, it is more meaningful to capture that in the reward function $\mathcal{F}_{P}$ rather than $A_{P}$.\par
 
% END Thmtheorem10.p2
 
% BEGIN Thmtheorem10.p3
 
A reward function $\mathcal{F}_{Q}$ for a reward interface $Q=(A_{Q},\mathcal{F}_{Q})$ is defined analogously. $\lrcorner$\par
 
% END Thmtheorem10.p3
 
% END Thmtheorem10
 
% BEGIN S3.p3
 
For $v\in\mathbb{R}_{-\infty}$ and $\sim\in\{<,\leq,\geq,>\}$, we define $\mathcal{F}_{P}^{\sim v}:=\{\sigma\in(\Sigma^{Obs}_{P})^{\infty}\mid\mathcal{F}_{P}(\sigma)\sim v\}$ to be the set of words which $\mathcal{F}_{P}$ maps to some value $\sim v$. When we write $\mathcal{F}_{P}(\sigma)\sim v$, we implicitly mean that $\mathcal{F}_{P}(\sigma)$ is defined. We define $\mathit{Vals}(\mathcal{F}_{P}):=\{v\in\mathbb{R}_{-\infty}\mid\exists\sigma\in(\Sigma^{Obs}_{P})^{\infty}.\nobreak\ \mathcal{F}(\sigma)=v\}$.\par
 
% END S3.p3
 
% BEGIN S3.p4
 
Let us fix an interface automaton $A=\langle S,S^{init},\Sigma^{I},\Sigma^{O},\Sigma^{E},\Sigma^{H},\mathcal{T}\rangle$ for the rest of this section. To evaluate implementations of an interface, we define the set of traces of $A$ as\par
  
\par  \quad  \[\begin{array}[]{lll}\mathsf{Traces}(A)&:=&\{\sigma\in\Sigma_{A}^{\omega}\mid\exists\nobreak\ \text{execution of }A\text{ on }\sigma\}\\
&\cup&\{\sigma\in\Sigma_{A}^{*}\mid\exists\nobreak\ \text{exec. of }A\text{ on }\sigma\text{ ending in }s:\mathsf{Enabled}(s,\Sigma^{O}\cup\Sigma^{H})=\emptyset\}.\end{array}\]\quad  \quad  
  
That is, $\mathsf{Traces}(A)$ is the set that consists of the infinite words over $\Sigma_{A}$ for which there exists an infinite execution, as well as the finite words over $\Sigma_{A}$ where an execution ends in a state with no output or internal action possible. That is, we consider maximal traces, taking into account that the environment inputs are not forced to occur.\par
 
% END S3.p4
 
% BEGIN S3.p5
 
For a given input sequence $\sigma_{E,I}\in(\Sigma^{E}\cup\Sigma^{I})^{\infty}$, we define $\mathsf{Traces}(A,\sigma_{E,I}):=\{\sigma\in\mathsf{Traces}(A)\mid\sigma|_{(\Sigma^{E}\cup\Sigma^{I})}=\sigma_{E,I}\}$ to be the subset of $\mathsf{Traces}(A)$ containing the words consistent with $\sigma_{E,I}$. These words represent all possible behaviors of $A$ when the environment provides the inputs specified by $\sigma_{E,I}$. If $\Gamma\subseteq\Sigma_{A}$ is an alphabet such that $\Gamma\supseteq\Sigma^{I}\cup\Sigma^{E}$, we define $\mathsf{Traces}_{\Gamma}(A,\sigma_{E,I}):=\{\gamma\in\Gamma^{\infty}\mid\exists\sigma\in\mathsf{Traces}(A,\sigma_{E,I})\land\sigma|_{\Gamma}=\gamma\}$ to be the projection of $\mathsf{Traces}(A,\sigma_{E,I})$ on $\Gamma$.\par
 
% END S3.p5
 
% BEGIN S3.p6
 
Reward functions impose no assumptions on the environment of a component. However, the value that a component can possibly achieve may depend on the behaviour of the environment. Therefore, we require that components implementing a reward interface satisfy the quantitative specification to the best extent possible with respect to the input provided by the component’s environment. This intuition is formalized by the good-enough criterion for interface automata (also used to model implementations).\par
 
% END S3.p6
 
% BEGIN S3.p7
 
First, we adapt to our setting the notion of hopeful sequences [2], which is used to characterize the input sequences for which a given reward value is possible. Note that in our setting, we consider asynchronous executions, such that there can be several consecutive inputs with no output in-between, or such that from some point on we only have outputs. This is in contrast to hopeful inputs in [2], which are defined for synchronous executions. We also allow arbitrary hidden actions, which may interleave with the non-internal actions.\par
 
% END S3.p7
 
% BEGIN Thmtheorem11
 
\par\medskip\noindent\textbf{Definition 11 (Hopeful Sequences).}\par
 
% BEGIN Thmtheorem11.p1
 
Given alphabets $\Sigma$ and $\Gamma$ such that $\Gamma\subseteq\Sigma$, a function $\mathcal{F}:\Sigma^{\infty}\to\mathbb{R}_{-\infty}$ and $v\in\mathbb{R}_{-\infty}$, we say that a sequence $\gamma\in\Gamma^{\infty}$ is $(\mathcal{F},v)$-hopeful if and only if there exists $\sigma\in\Sigma^{\infty}$ such that $\gamma=\sigma|_{\Gamma}$, $\mathcal{F}(\sigma)$ is defined, and $\mathcal{F}(\sigma)\geq v$. We denote the set of $(\mathcal{F},v)$-hopeful $\Gamma$-sequences by $\mathit{Hopeful}(\mathcal{F},v,\Gamma):=\{\gamma\in\Gamma^{\infty}\mid\exists\sigma\in\Sigma^{\infty}.\nobreak\ \gamma=\sigma|_{\Gamma}\land\mathcal{F}(\sigma)\geq v\}.$\par
 
% END Thmtheorem11.p1
 
% END Thmtheorem11
 
% BEGIN S3.p8
 
We use $\Gamma$ here to denote a subset of $\Sigma$, which will typically be instantiated to be a subset of the inputs $\Sigma^{I}\cup\Sigma^{E}$, as we will see in the following definition. This notion describes the potential quality of the environment, characterizing the inputs to an interface as hopeful with respect only to specific values of $\mathcal{F}$. The hopefulness of an input sequence does not depend on the interface automaton itself, but only on the reward function $\mathcal{F}$. An interface automaton is good-enough with respect to a reward function $\mathcal{F}$, if, intuitively, the traces it generates on any $(\mathcal{F},v)$-hopeful input sequence achieve reward at least $v$.\par
 
% END S3.p8
 
% BEGIN Thmtheorem12
 
\par\medskip\noindent\textbf{Definition 12 (Good-Enough Interface Automaton).}\par
 
% BEGIN Thmtheorem12.p1
 
Consider an interface automaton $A$ and a reward function $\mathcal{F}:\Delta^{\infty}\to\mathbb{R}_{-\infty}$. We say that $A$ is good-enough with respect to $\mathcal{F}$ if and only if for every value $v\in\mathit{Vals}(\mathcal{F})$, every input sequence $\sigma_{E,I}\in\mathit{Hopeful}(\mathcal{F},v,(\Sigma_{A}^{E}\cup\Sigma_{A}^{I})\cap\Delta)$, and every $\sigma\in\mathsf{Traces}_{(\Delta\cap\Sigma_{A})}(A,\sigma_{E,I})$ it holds that $\mathcal{F}(\sigma)$ is defined and $\mathcal{F}(\sigma)\geq v$.\par
 
% END Thmtheorem12.p1
 
% END Thmtheorem12
 
% BEGIN S3.p9
 
This definition is in the spirit of [2], as a good-enough interface automaton must perform to the best extent possible according to $\mathcal{F}$, but only for the input it receives. In Definition 12 we used a general alphabet $\Delta$. This is useful for the definition of compatibility of reward interfaces with respect to a given reward function.\par
 
% END S3.p9
 
% BEGIN S3.p10
 
Now we define the notion of (best-effort) implementation of a reward interface, which must be good-enough with respect the interface’s reward function.\par
 
% END S3.p10
 
% BEGIN Thmtheorem13
 
\par\medskip\noindent\textbf{Definition 13 (Best-Effort Implementation).}\par
 
% BEGIN Thmtheorem13.p1
 
An interface automaton $S$ is an implementation of a reward interface $P=(A_{P},\mathcal{F}_{P})$ if and only if it satisfies the following conditions.\par
  
\par\noindent  1. 
% BEGIN S3.I1.i1.p1
 
$S$ refines the interface automaton $A_{P}$.\par
 
% END S3.I1.i1.p1
 
 
\par\noindent  2. 
% BEGIN S3.I1.i2.p1
 
$S$ is good-enough with respect to the function $\mathcal{F}_{P}$.\par
 
% END S3.I1.i2.p1
 
  
We denote with $\mathsf{Imp}(P)$ the set of all implementations of $P$.\par
 
% END Thmtheorem13.p1
 
% END Thmtheorem13
 
% BEGIN S3.p11
 
An implementation of a reward interface $P$ is an instance of a classical interface automaton, as it does not have a reward function. In the original theory, there is no distinction between an interface and its implementation, as the interface is defined by the automaton structure alone. When considering an interface in isolation, the external inputs $\Sigma^{E}$ and inputs $\Sigma^{I}$ are treated in the same way. However, we differentiate between the two when considering the interface in the context of rest of the system, which generates the inputs $\Sigma^{I}$.\par
 
% END S3.p11
 
% BEGIN Thmtheorem14
 
\par\medskip\noindent\textbf{Example 14.}\par
 
% BEGIN Thmtheorem14.p1
 
Let us examine how a best-effort implementation for $P$ could act. With respect to the reward function $\mathcal{F}_{P}$ from Example 10, a possible implementation $S_{P}$ for $P$ waits for $\mathsf{e_{1}}$ or $\mathsf{e_{2}}$, and then expects input $\mathsf{q}$ to follow with $\mathsf{o_{1}}$, or produces output $\mathsf{p}$, respectively. We give a portion of the respective implementation in Figure 2(a), for the behaviors where $\mathsf{e_{1}}$ occurs first. The missing part has analogous structure for the case when $\mathsf{e_{2}}$ is received first.\par
 
% END Thmtheorem14.p1
 
% BEGIN Thmtheorem14.p2
 
As the implementation must not perform $\mathsf{o_{1}}$ unless $\mathsf{e_{1}}$ was received, it waits until $\mathsf{e_{1}}$ happens. Should $\mathsf{q}$ never appear, $S_{P}$ has no obligation to perform $\mathsf{o_{1}}$. This is because, even though $\mathcal{F}_{P}$ does not require $\mathsf{q}$, Definition 12 considers the traces of $S_{P}$, where $\mathsf{q}$ must be read before $\mathsf{o_{1}}$ can be produced. While $P$ is unable to control or observe $\mathsf{o_{2}}$, it exercises control by withholding $\mathsf{p}$ until $\mathsf{e_{2}}$, which is represented in $\mathcal{F}_{P}$, and therefore in implementation $S_{P}$. If instead, the implementation could immediately produce $\mathsf{p}$, the value for $\mathcal{F}_{P}$ would be lower since it could be that no $\mathsf{e_{2}}$ occurred. $\lrcorner$\par
 
% END Thmtheorem14.p2
 
% BEGIN S3.F2
   
% BEGIN S3.F2.sf1

\begin{center}\includegraphics[width=0.95\linewidth,height=0.45\textheight,keepaspectratio]{graphic-03.pdf}\end{center}
 
(a) (Part of) implementation $S_{P}$ of $P$.\par
 
% END S3.F2.sf1
   
% BEGIN S3.F2.sf2

\begin{center}\includegraphics[width=0.95\linewidth,height=0.45\textheight,keepaspectratio]{graphic-04.pdf}\end{center}
 
(b) Composition $A_{P}\parallel A_{Q}$.\par
 
% END S3.F2.sf2
   
Figure 2: Interface automata for an implementation of $P$ and composition of $A_{P}$ and $A_{Q}$.\par
 
% END S3.F2
 
% END Thmtheorem14
 
% END S3


% BEGIN S4
 
\par\medskip\noindent\textbf{4 Compatibility and Composition of Reward Interfaces}\par
 
% BEGIN S4.p1
 
In this section, we lift the notions of interface compatibility and composition to reward interfaces. Following the classical interface theory, compatibility of reward interfaces requires the existence of a legal environment for the underlying interface automata. For the quantitative part of our reward interfaces, we define compatibility with respect to a “joint” reward function $\mathcal{F}$, which, intuitively, is a quantitative specification with respect to which the composition of the implementations of the two interfaces must be good enough. In top-down component-based design, we need to ensure that the local reward functions for interfaces guarantee that the composition of their implementations is good-enough with respect to a given high-level reward function $\mathcal{F}$. This is precisely the $\mathcal{F}$-compatibility criterion we define.\par
 
% END S4.p1
 
% BEGIN S4.p2
 
For this section, let $P=(A_{P},\mathcal{F}_{P})$ and $Q=(A_{Q},\mathcal{F}_{Q})$ be composable reward interfaces with $A_{P}=\langle V_{P},V^{init}_{P},\Sigma^{I}_{P},\Sigma^{O}_{P},\Sigma^{E}_{P},\Sigma^{H}_{P},\mathcal{T}_{P}\rangle$ and $A_{Q}=\langle V_{Q},V^{init}_{Q},\Sigma^{I}_{Q},\Sigma^{O}_{Q},\Sigma^{E}_{Q},\Sigma^{H}_{Q},\mathcal{T}_{Q}\rangle$, and $A_{P}\otimes A_{Q}=\langle V_{P\otimes Q},V_{P\otimes Q}^{init},\Sigma_{P\otimes Q}^{I},\Sigma_{P\otimes Q}^{O},\Sigma_{P\otimes Q}^{E},\Sigma_{P\otimes Q}^{H},\mathcal{T}_{P\otimes Q}\rangle$ their product [Definition 4].\par
 
% END S4.p2
 
% BEGIN S4.p3
 
We now define $\mathcal{F}$-compatibility, for a given reward function $\mathcal{F}:(\Sigma_{P\otimes Q}^{Obs})^{\infty}\to\mathbb{R}_{-\infty}$ with alphabet consisting of the non-internal actions of the product automaton $A_{P}\otimes A_{Q}$.\par
 
% END S4.p3
 
% BEGIN Thmtheorem15
 
\par\medskip\noindent\textbf{Definition 15 ($\mathcal{F}$-Compatibility).}\par
 
% BEGIN Thmtheorem15.p1
 
The reward interfaces $P=(A_{P},\mathcal{F}_{P})$ and $Q=(A_{Q},\mathcal{F}_{Q})$ are $\mathcal{F}$-compatible for a given function $\mathcal{F}:(\Sigma_{P\otimes Q}^{Obs})^{\infty}\to\mathbb{R}_{-\infty}$ if and only if $A_{P}$ and $A_{Q}$ are compatible and for all implementations $S_{P}\in\mathsf{Imp}(P)$ and $S_{Q}\in\mathsf{Imp}(Q)$ of $P$ and $Q$ respectively, if $S_{P}$ and $S_{Q}$ are composable, then $S_{P}\parallel S_{Q}$ is good-enough with respect to $\mathcal{F}$.\par
 
% END Thmtheorem15.p1
 
% END Thmtheorem15
 
% BEGIN Thmtheorem16
 
\par\medskip\noindent\textbf{Example 16.}\par
 
% BEGIN Thmtheorem16.p1
 
We now illustrate that the reward interfaces $P$ and $Q$ are $\mathcal{F}_{\mathsf{S}}$-compatible for the reward function $\mathcal{F}_{\mathsf{S}}$ from Example 1. In particular, the composition of $P$ and $Q$ refines the high-level specification of $\mathsf{S}$ expressed by $\mathcal{F}_{\mathsf{S}}$. Clearly, $P$ and $Q$ are composable, and since $\mathsf{Illegal}(A_{P},A_{Q})=\emptyset$, we can construct their composition $A_{P}\parallel A_{Q}$ (shown in Figure 2(b)).\par
 
% END Thmtheorem16.p1
 
% BEGIN Thmtheorem16.p2
 
To see that $P$ and $Q$ are $\mathcal{F}_{\mathsf{S}}$-compatible, consider any pair of implementations $S_{P}$ and $S_{Q}$. Now we explain why the composition $S_{P}\parallel S_{Q}$ must be good-enough with respect to $\mathcal{F}_{\mathsf{S}}$. We examine $(\mathcal{F}_{\mathsf{S}},v)$-hopeful input sequences, and consider the expected behaviors of $S_{P}$ and $S_{Q}$. $\mathcal{F}_{P}$ requires $S_{P}$ to produce $\mathsf{o_{1}}$ once $\mathsf{e_{1}}$ was received, and not before. The same holds with respect to $\mathcal{F}_{Q}$, and $\mathsf{o_{2}}$ and $\mathsf{e_{2}}$. Since $\mathcal{F}_{P}$ and $\mathcal{F}_{Q}$ further constrain $\mathsf{p}$ and $\mathsf{q}$, respectively, this ensures that both $S_{P}$ and $S_{Q}$ will each have the opportunity to output $\mathsf{o_{1}}$ or $\mathsf{o_{2}}$ respectively, once required. Thus, any pair of best-effort implementations of $P$ and $Q$ together will operate as a best-effort implementation of the high-level reward function $\mathcal{F}_{\mathsf{S}}$.\par
 
% END Thmtheorem16.p2
 
% BEGIN Thmtheorem16.p3
 
In top-down design, we may further refine our model, by splitting $Q$ into two components $Q_{1}$ and $Q_{2}$. Suppose for example that $Q_{1}$ handles $\mathsf{p}$, and can output $\mathsf{o_{3}}$ and some $\mathsf{q}_{1}$. $Q_{2}$ outputs $\mathsf{o_{2}}$ and $\mathsf{q}$, and synchronizes with $Q_{1}$ via $\mathsf{q}_{1}$. The interface automata $A_{Q_{1}}$ and $A_{Q_{2}}$ are in Figure 3. The reward function $\mathcal{F}_{Q_{1}}$ requires that $\mathsf{q}_{1}$ only occurs once $\mathsf{e_{2}}$ did. Similarly, $\mathcal{F}_{Q_{2}}$ requires $\mathsf{q}$ after $\mathsf{e_{1}}$. Through action $\mathsf{q}_{1}$, $Q_{1}$ can prevent $Q_{2}$ from incorrectly producing $\mathsf{o_{2}}$.\par
 
% END Thmtheorem16.p3
 
% BEGIN Thmtheorem16.p4
 
In the product $A_{Q_{1}}\otimes A_{Q_{2}}$, the state $(v_{1},w_{2})$ is illegal, because $A_{Q_{1}}$ can output $\mathsf{q}_{1}$ from $v_{1}$, but $w_{2}$ in $A_{Q_{2}}$ can not accept it. In the composition $A_{Q_{1}}\parallel A_{Q_{2}}$, state $(v_{1},w_{2})$ is therefore removed. It is easy to see that the composition of any pair of implementations of $Q_{1}$ and $Q_{2}$ is good enough with respect to $\mathcal{F}_{Q}$. Thus, $Q_{1}$ and $Q_{2}$ are $\mathcal{F}_{Q}$-compatible. $\lrcorner$\par
 
% END Thmtheorem16.p4
 
% BEGIN S4.F3
   
% BEGIN S4.F3.sf1

\begin{center}\includegraphics[width=0.95\linewidth,height=0.45\textheight,keepaspectratio]{graphic-05.pdf}\end{center}
 
(a) Interface automaton $A_{Q_{1}}$.\par
 
% END S4.F3.sf1
   
% BEGIN S4.F3.sf2

\begin{center}\includegraphics[width=0.95\linewidth,height=0.45\textheight,keepaspectratio]{graphic-06.pdf}\end{center}
 
(b) Interface automaton $A_{Q_{2}}$.\par
 
% END S4.F3.sf2
   
Figure 3: Splitting of $Q$ into two components $Q_{1}$ and $Q_{2}$ as described in Example 16.\par
 
% END S4.F3
 
% END Thmtheorem16
 
% BEGIN S4.p4
 
The $\mathcal{F}$-compatibility of $A_{P}$ and $A_{Q}$ guarantees that $(A_{P}\parallel A_{Q},\mathcal{F})$ is a reward interface, as it requires that the interface automata $A_{P}$ and $A_{Q}$ are compatible, and $\mathcal{F}$ has the right domain. Furthermore, the second condition of $\mathcal{F}$-compatibility ensures that $P$ and $Q$ can be implemented independently, and composing the resulting implementations yields a good-enough implementation of $(A_{P}\parallel A_{Q},\mathcal{F})$. By checking $\mathcal{F}$-compatibility of two reward interfaces for a given reward function $\mathcal{F}$, we can establish whether their individual reward functions are aligned to guarantee the good-enough satisfaction of the specification $\mathcal{F}$.\par
 
% END S4.p4
 
% BEGIN S4.p5
 
In bottom-up design, on the other hand, we want to construct the composition of two interfaces by combining their reward functions, in order to express the guarantees of two interfaces when put together. For this, we need to provide a definition of composition of reward functions. The definition is parametrized by a function $\mathit{comb}:\mathbb{R}_{-\infty}\times\mathbb{R}_{-\infty}\to\mathbb{R}_{-\infty}$, which specifies how we combine the numeric values assigned by the two functions.\par
 
% END S4.p5
 
% BEGIN S4.p6
 
The observable behaviours of the composition $P\parallel Q$ are over the alphabet $\Sigma_{P\otimes Q}^{Obs}=(\Sigma_{P}^{Obs}\cup\Sigma_{Q}^{Obs})\setminus\mathsf{Shared}(P,Q)$ of non-internal actions of $P\parallel Q$. Thus, the domain of the composed reward function must be the set of sequences over $\Sigma_{P\otimes Q}^{Obs}$. Our goal is to define the composition of $\mathcal{F}_{P}$ and $\mathcal{F}_{Q}$ in a way that captures the combined value achieved by any pair of implementations of the two interfaces. Since the implementations of each of the interfaces are constrained by the respective reward function, we cannot assume that any pair of implementations will be cooperating. Therefore, the value assigned to a sequence $\sigma\in(\Sigma_{P\otimes Q}^{Obs})^{\infty}$ corresponds to the worst case over all the traces produced by some pair of composable implementations that are compatible with the sequence $\sigma|_{\Sigma_{P\otimes Q}^{I,E}}$ of inputs in $\sigma$.\par
 
% END S4.p6
 
% BEGIN Thmtheorem17
 
\par\medskip\noindent\textbf{Definition 17 (Reward Function Composition).}\par
 
% BEGIN Thmtheorem17.p1
 
Let $\mathit{comb}:\mathbb{R}_{-\infty}\times\mathbb{R}_{-\infty}\to\mathbb{R}_{-\infty}$. We define the composition ${{\mathcal{F}_{P}}\mathbin{\triangledown_{\mathit{comb}}}{\mathcal{F}_{Q}}}:(\Sigma_{P\otimes Q}^{Obs})^{\infty}\to\mathbb{R}_{-\infty}$ of the functions $\mathcal{F}_{P}$ and $\mathcal{F}_{Q}$, such that for $\sigma\in(\Sigma_{P\otimes Q}^{Obs})^{\infty}$ where there exist composable $S_{P}^{\prime}\in\mathsf{Imp}(P)$ and $S_{Q}^{\prime}\in\mathsf{Imp}(Q)$ and $\sigma^{\prime}\in\mathsf{Traces}(S_{P}^{\prime}\parallel S_{Q}^{\prime})$ such that $\sigma^{\prime}|_{\Sigma_{P\otimes Q}^{Obs}}=\sigma$, we let\par
  
\par  \quad  \[\begin{array}[]{ll}{{\mathcal{F}_{P}}\mathbin{\triangledown_{\mathit{comb}}}{\mathcal{F}_{Q}}}(\sigma):=\inf\{&\mathit{comb}(\mathcal{F}_{P}(\sigma^{\prime\prime}|_{\Sigma_{P}^{Obs}}),\mathcal{F}_{Q}(\sigma^{\prime\prime}|_{\Sigma_{Q}^{Obs}}))\mid S_{P}\in\mathsf{Imp}(P)\text{ and }\\
&S_{Q}\in\mathsf{Imp}(Q)\text{ composable},\sigma^{\prime\prime}\in\mathsf{Traces}(S_{P}\parallel S_{Q},\sigma|_{\Sigma_{P\otimes Q}^{I,E}})\},\end{array}\]\quad  \quad  
  
and ${{\mathcal{F}_{P}}\mathbin{\triangledown_{\mathit{comb}}}{\mathcal{F}_{Q}}}(\sigma)$ is undefined otherwise.\par
 
% END Thmtheorem17.p1
 
% END Thmtheorem17
 
% BEGIN S4.p7
 
By taking the worst case in the definition of ${{\mathcal{F}_{P}}\mathbin{\triangledown_{\mathit{comb}}}{\mathcal{F}_{Q}}}$, we guarantee that as long as $A_{P}$ and $A_{Q}$ are compatible, the reward interfaces $P$ and $Q$ are ${{\mathcal{F}_{P}}\mathbin{\triangledown_{\mathit{comb}}}{\mathcal{F}_{Q}}}$-compatible, as shown in the next proposition. Composing $P$ and $Q$ results in the reward interface $(A_{P}\parallel A_{Q},{{\mathcal{F}_{P}}\mathbin{\triangledown_{\mathit{comb}}}{\mathcal{F}_{Q}}})$. $P$ and $Q$ can be implemented independently, and the composition of their implementations will be good enough with respect to ${{\mathcal{F}_{P}}\mathbin{\triangledown_{\mathit{comb}}}{\mathcal{F}_{Q}}}$.\par
 
% END S4.p7
 
% BEGIN Thmtheorem18
 
\par\medskip\noindent\textbf{Proposition 18 (${{\mathcal{F}_{P}}\mathbin{\triangledown_{\mathit{comb}}}{\mathcal{F}_{Q}}}$-Compatibility).}\par
 
% BEGIN Thmtheorem18.p1
 
Let $\mathit{comb}:\mathbb{R}_{-\infty}\times\mathbb{R}_{-\infty}\to\mathbb{R}_{-\infty}$. For all reward interfaces $P=(A_{P},\mathcal{F}_{P})$ and $Q=(A_{Q},\mathcal{F}_{Q})$ with $A_{P}$ and $A_{Q}$ compatible, it holds that $P$ and $Q$ are ${{\mathcal{F}_{P}}\mathbin{\triangledown_{\mathit{comb}}}{\mathcal{F}_{Q}}}$-compatible.\par
 
% END Thmtheorem18.p1
 
% END Thmtheorem18
 
% BEGIN S4.p8
 
We discard any observable traces not resulting from implementations of $P$ and $Q$, restricting the observable behaviors of interface automata good enough w.r.t. ${{\mathcal{F}_{P}}\mathbin{\triangledown_{\mathit{comb}}}{\mathcal{F}_{Q}}}$ to those resulting from some $S_{P}\parallel S_{Q}$ where $S_{P}\in\mathsf{Imp}(P)$ and $S_{Q}\in\mathsf{Imp}(Q)$. This composition is associative if the function $\mathit{comb}$ used to combine the values is associative, and satisfies some monotonicity and continuity conditions (Proposition 29).\par
 
% END S4.p8
 
% BEGIN Thmtheorem19
 
\par\medskip\noindent\textbf{Proposition 19 (Quality of the Reward Function Composition).}\par
 
% BEGIN Thmtheorem19.p1
 
Consider reward interfaces $P=(A_{P},\mathcal{F}_{P})$ and $Q=(A_{Q},\mathcal{F}_{Q})$ with $A_{P}$ and $A_{Q}$ compatible, and a reward function $\mathcal{F}$ such that $P$ and $Q$ are $\mathcal{F}$-compatible. Then, $\mathsf{Imp}(A_{P}\parallel A_{Q},{{\mathcal{F}_{P}}\mathbin{\triangledown_{\mathit{comb}}}{\mathcal{F}_{Q}}})\subseteq\mathsf{Imp}(A_{P}\parallel A_{Q},\mathcal{F})$.\par
 
% END Thmtheorem19.p1
 
% END Thmtheorem19
 
% BEGIN Thmtheorem20
 
\par\medskip\noindent\textbf{Example 20.}\par
 
% BEGIN Thmtheorem20.p1
 
Let us consider components $P$ and $Q$ from Example 1 from the perspective of bottom-up design. Recall the reward function $\mathcal{F}_{P}$ from Example 14, where we impose that for input sequences containing both $\mathsf{e_{1}}$ and $\mathsf{e_{2}}$, there must be outputs $\mathsf{o_{1}}$ and $\mathsf{p}$ in the respective order. Now, suppose instead we use a function $\mathcal{F}_{P}^{\prime}$, which ignores the requirements on $\mathsf{p}$, meaning that it only asks for $\mathsf{o_{1}}$ to occur after $\mathsf{e_{1}}$ and $\mathsf{q}$, and nothing else. For those input sequences with both $\mathsf{e_{1}}$ and $\mathsf{e_{2}}$, $\mathcal{F}_{P}^{\prime}$ will always award a value of $\frac{1}{2}$ as long as $\mathsf{o_{1}}$ is produced. Note that this does not disallow $\mathsf{p}$, it simply does not require it.\par
 
% END Thmtheorem20.p1
 
% BEGIN Thmtheorem20.p2
 
A possible behavior of an implementation $S_{P}^{\prime}$ of $(A_{P},\mathcal{F}_{P}^{\prime})$ is then to produce $\mathsf{o_{1}}$ once receiving $\mathsf{q}$, and be idle otherwise. In particular, it will never produce $\mathsf{p}$. Implementations $S_{Q}$ of $Q$ will, as before, produce $\mathsf{o_{2}}$ only after receiving $\mathsf{e_{2}}$ and $\mathsf{p}$.\par
 
% END Thmtheorem20.p2
 
% BEGIN Thmtheorem20.p3
 
It is easy to see that the new function $\mathcal{F}_{P}^{\prime}$ allows implementations of $(A_{P},\mathcal{F}_{P}^{\prime})$ which prevent $Q$ from producing $\mathsf{o_{2}}$, disabling the best value w.r.t. $\mathcal{F}_{Q}$. Since $P$ has no knowledge of $\mathcal{F}_{Q}$, we cannot assume that $S_{P}$ will be more cooperative than required by $A_{P}$ and $\mathcal{F}_{P}^{\prime}$. As $A_{P}$ and $A_{Q}$ are compatible, choosing $\mathit{comb}$ to be the sum, we have the guarantee that any $S_{P}\parallel S_{Q}$ is good enough w.r.t. ${{\mathcal{F}_{P}^{\prime}}\mathbin{\triangledown_{\mathit{comb}}}{\mathcal{F}_{Q}}}$. This guarantee is however weaker than $\mathcal{F}_{\mathsf{S}}$. In this case, that means that for the mentioned input sequences featuring $\mathsf{e_{1}}$ and $\mathsf{e_{2}}$, where previously the best value of $\frac{1}{2}+\frac{1}{2}$ was possible, we will now achieve at least $\frac{1}{2}+0$.\par
 
% END Thmtheorem20.p3
 
% BEGIN Thmtheorem20.p4
 
In bottom-up design, we can use $(A_{P}\parallel A_{Q},{{\mathcal{F}_{P}^{\prime}}\mathbin{\triangledown_{\mathit{comb}}}{\mathcal{F}_{Q}}})$ as the reward interface for the composition of $P^{\prime}=(A_{P},\mathcal{F}_{P}^{\prime})$ and $Q$. This captures the guarantees of the composed components at a higher level of the system, without the need to devise a new reward function and checking if $P$ and $Q$ are compatible with respect to it. $\lrcorner$\par
 
% END Thmtheorem20.p4
 
% END Thmtheorem20
 
% END S4


% BEGIN S5
 
\par\medskip\noindent\textbf{5 Reward Interface Refinement}\par
 
% BEGIN S5.p1
 
Supporting different levels of abstraction for the same component is helpful in the design process, as it allows a coarse description of interactions and dependencies while also defining internal details. Refinement enables this independent implementability of an interface, permitting extended functionality as long as original obligations are maintained. For a reward interface $Q$ to be a refinement of a reward interface $P$, we require as usual that $Q$ works within the same environments as $P$, and that an implementation of $Q$ is also an implementation of $P$. This means, in particular, that, similarly to classical interface automata, a refinement of $P$ must accept at least all inputs of $P$ and not produce more outputs than $P$. In addition, we must consider the reward function w.r.t. best-effort implementations.\par
 
% END S5.p1
 
% BEGIN S5.p2
 
The refinement relation for reward interfaces can be defined as the conjunction of refinement between the respective interface automata [Definition 8] and refinement between the respective reward functions. We define the notion of refinement for quantitative reward functions in the context of good-enough satisfaction. Intuitively, for a function $\mathcal{F}_{Q}$ to refine a function $\mathcal{F}_{P}$, each input sequence for $Q$ must be “as hopeful” w.r.t. $\mathcal{F}_{Q}$ as it is w.r.t. $\mathcal{F}_{P}$, and the value assigned by $\mathcal{F}_{Q}$ to a trace must be matched by the value assigned by $\mathcal{F}_{P}$ to the corresponding trace in $P$. This captures the idea that $\mathcal{F}_{Q}$ inherits the obligations of $\mathcal{F}_{P}$ and assigns values in a way that is possible in $\mathcal{F}_{P}$. We formalize this intuition as the existence of a helper function $r_{PQ}:\mathbb{R}_{-\infty}\to\mathbb{R}_{-\infty}$, that suitably relates the ranges $\mathit{Vals}(\mathcal{F}_{P})$ and $\mathit{Vals}(\mathcal{F}_{Q})$ of the functions. The idea is that $r_{PQ}$ preserves the relative order of the values in $\mathcal{F}_{P}$, as it relates to both the hopefulness of input traces and values assigned to the respective full traces.\par
 
% END S5.p2
 
% BEGIN Thmtheorem21
 
\par\medskip\noindent\textbf{Definition 21 (Reward Function Refinement).}\par
 
% BEGIN Thmtheorem21.p1
 
Suppose that for the reward interfaces $P=(A_{P},\mathcal{F}_{P})$ and $Q=(A_{Q},\mathcal{F}_{Q})$ we have $\Sigma_{P}^{I}\subseteq\Sigma_{Q}^{I}$, $\Sigma_{P}^{E}\subseteq\Sigma_{Q}^{E}$ and $\Sigma_{Q}^{O}\subseteq\Sigma_{P}^{O}$. We say that $\mathcal{F}_{Q}$ refines $\mathcal{F}_{P}$, denoted $\mathcal{F}_{Q}\preceq\mathcal{F}_{P}$, iff there exists a function $r_{PQ}:\mathbb{R}_{-\infty}\to\mathbb{R}_{-\infty}$ such that for every value $v\in\mathit{Vals}(\mathcal{F}_{P})$, the following two conditions are satisfied.\par
  
\par\noindent  1. 
% BEGIN S5.I1.i1.p1
 
For every $\sigma_{Q}^{I,E}\in(\Sigma_{Q}^{I}\cup\Sigma_{Q}^{E})^{\infty}$, if $\sigma_{Q}^{I,E}|_{\Sigma_{P}}\in\mathit{Hopeful}(\mathcal{F}_{P},v,\Sigma_{P}^{I}\cup\Sigma_{P}^{E})$, For every $\sigma_{Q}^{I,E}\in(\Sigma_{Q}^{I}\cup\Sigma_{Q}^{E})^{\infty}$, then $\sigma_{Q}^{I,E}\in\mathit{Hopeful}(\mathcal{F}_{Q},r_{PQ}(v),\Sigma_{Q}^{I}\cup\Sigma_{Q}^{E})$.\par
 
% END S5.I1.i1.p1
 
 
\par\noindent  2. 
% BEGIN S5.I1.i2.p1
 
For every $\sigma_{Q}\in(\Sigma_{Q}^{Obs})^{\infty}$, if $\mathcal{F}_{Q}(\sigma_{Q})\geq r_{PQ}(v)$ then $\mathcal{F}_{P}(\sigma_{Q}|_{\Sigma_{P}})\geq v$.\par
 
% END S5.I1.i2.p1
 
  
% END Thmtheorem21.p1
 
% END Thmtheorem21
 
% BEGIN S5.p3
 
Definition 21 captures the nature of refinement, because while it retains priority brackets for obligations of the original function, we can also introduce new dependencies or arbitrary subdivisions within each priority. The next example illustrates the refinement relation for reward functions and demonstrates the use of the helper function for establishing the relation.\par
 
% END S5.p3
 
% BEGIN Thmtheorem22
 
\par\medskip\noindent\textbf{Example 22.}\par
 
% BEGIN Thmtheorem22.p1
 
Going back to the reward interface for $\mathsf{S}$ with $\mathcal{F}_{\mathsf{S}}$ from Example 1, suppose we want to modify the guarantees, requiring $\mathsf{S}$ to output $\mathsf{o_{1}}$ or $\mathsf{o_{2}}$ for a period of time. More concretely, we ask for $\mathsf{o_{1}}$ to hold as often as possible from the point $\mathsf{e_{1}}$ is first received until the first occurrence of $\mathsf{e_{2}}$, as to capture the change in input. We add the same requirement on $\mathsf{o_{2}}$ if $\mathsf{e_{2}}$ occurs first. We modify the existing reward function $\mathcal{F}_{\mathsf{S}}$ to $\mathcal{F}_{R}$, where we multiply the assigned value by a factor of $\frac{2}{3}$ if this new condition is violated.\par
 
% END Thmtheorem22.p1
 
% BEGIN Thmtheorem22.p2
 
Clearly, a subset of traces that had value $1$ for $\mathcal{F}_{\mathsf{S}}$ are now assigned value $\frac{2}{3}$, namely where the other output $\mathsf{o_{3}}$ is interjected. The same holds for traces going from $\frac{1}{2}$ to $\frac{1}{3}$. The new requirement will always be violated if the order was not preserved in the first place, giving us $\frac{1}{6}$ instead of $\frac{1}{4}$. However, the old obligations are kept: Any $(\mathcal{F}_{\mathsf{S}},1)$-hopeful input will be paired with a trace that for $\mathcal{F}_{R}$ achieves at least $\frac{2}{3}$. The function $r_{\mathsf{S}R}$ will then be $r_{\mathsf{S}R}=\{(0,0),(\frac{1}{4},\frac{1}{6}),(\frac{1}{2},\frac{1}{3}),(1,\frac{2}{3})\}$. Thus, the priorities from $\mathcal{F}_{\mathsf{S}}$ are preserved in the refinement, even though we are able to strengthen requirements in the reward function $\mathcal{F}_{R}$. $\lrcorner$\par
 
% END Thmtheorem22.p2
 
% END Thmtheorem22
 
% BEGIN S5.p4
 
The refinement relation for reward interfaces combines the two refinement relations.\par
 
% END S5.p4
 
% BEGIN Thmtheorem23
 
\par\medskip\noindent\textbf{Definition 23 (Reward Interface Refinement).}\par
 
% BEGIN Thmtheorem23.p1
 
A reward interface $Q=(A_{Q},\mathcal{F}_{Q})$ refines a reward interface $P=(A_{P},\mathcal{F}_{P})$, written as $Q\preceq P$, if and only if $A_{Q}\preceq A_{P}$ and $\mathcal{F}_{Q}\preceq\mathcal{F}_{P}$.\par
 
% END Thmtheorem23.p1
 
% END Thmtheorem23
 
% BEGIN S5.p5
 
The rest of this section is dedicated to establishing that the refinement relation we defined has the properties necessary for compositional design. We establish that the refinement relation $\preceq$ between reward interfaces is a preorder, i.e. it is reflexive and transitive (Proposition 30). In particular, transitivity of refinement is important, as it allows for an iterative design process, where each refinement needs to only consider the last to maintain a refinement relation to the original specification. It is established by composing the respective helper functions.\par
 
% END S5.p5
 
% BEGIN S5.p6
 
As a consequence, an implementation of a refinement is then also an implementation of the original interface. Since the implementation must be good-enough w.r.t. the refined function, by refinement we can show it is also good-enough w.r.t. the original function.\par
 
% END S5.p6
 
% BEGIN Thmtheorem24
 
\par\medskip\noindent\textbf{Theorem 24 (Implementation of a Refinement).}\par
 
% BEGIN Thmtheorem24.p1
 
If a reward interface $Q=(A_{Q},\mathcal{F}_{Q})$ refines a reward interface $P=(A_{P},\mathcal{F}_{P})$, then $\mathsf{Imp}(Q)\subseteq\mathsf{Imp}(P)$.\par
 
% END Thmtheorem24.p1
 
% END Thmtheorem24
 
% BEGIN S5.p7
 
The other direction of Theorem 24 does not hold in general. First of all, a reward interface constrains the set of implementations through both the interface automaton and the reward function, while the refinement relation treats those two components separately, which makes it stronger than the inclusion relation between the sets of implementations. Moreover, the refinement relation on reward functions alone is stronger than the inclusion between the respective sets of good-enough automata.\par
 
% END S5.p7
 
% BEGIN S5.p8
 
The last properties we establish enable independent implementability, by guaranteeing that refinement does not impair the original context of an interface. They concern refinement in the function composition and the property of substitutability. Intuitively, substitutability states that a refinement of a reward interface $P$ remains $\mathcal{F}$-compatible with any reward interfaces that are $\mathcal{F}$-compatible with $P$.\par
 
% END S5.p8
 
% BEGIN Thmtheorem25
 
\par\medskip\noindent\textbf{Theorem 25.}\par
 
% BEGIN Thmtheorem25.p1
 
Consider $P=(A_{P},\mathcal{F}_{P})$ and $Q=(A_{Q},\mathcal{F}_{Q})$ with $A_{P}$ and $A_{Q}$ compatible, and let $P^{\prime}=(A_{P^{\prime}},\mathcal{F}_{P^{\prime}})$ be a non-empty refinement of $P$, with $(\Sigma_{P^{\prime}}^{I}\setminus\Sigma_{P}^{I})\cap\Sigma_{Q}^{O}=\emptyset$. Then if $\mathit{comb}$ is monotonically increasing, $(A_{P^{\prime}}\parallel A_{Q},{{\mathcal{F}_{P^{\prime}}}\mathbin{\triangledown_{\mathit{comb}}}{\mathcal{F}_{Q}}})\preceq(A_{P}\parallel A_{Q},{{\mathcal{F}_{P}}\mathbin{\triangledown_{\mathit{comb}}}{\mathcal{F}_{Q}}})$. If $P$ and $Q$ are $\mathcal{F}$-compatible for a reward function $\mathcal{F}$, then $P^{\prime}$ and $Q$ are also $\mathcal{F}$-compatible.\par
 
% END Thmtheorem25.p1
 
% END Thmtheorem25
 
% BEGIN S5.p9
 
With that, we established that the refinement for reward interfaces in Definition 23 indeed lifts the ideas of interface refinement to the best-effort quantitative setting and has the properties to enable component-based design of systems with quantitative specifications.\par
 
% END S5.p9
 
% END S5


% BEGIN S6
 
\par\medskip\noindent\textbf{6 Checking Compatibility, Refinement, and Implementability}\par
 
% BEGIN S6.p1
 
In this section, we study the algorithmic aspects of our theory of reward interfaces, in the context of an automata-based finite representation of the reward functions. We first introduce this representation. Then, we define the decision problems of interest, namely checking compatibility, refinement and implementability of reward interfaces, and outline algorithms for each of them for our representation. Lastly, we discuss quantitative temporal logics and weighted automata as representations of reward functions, and relate them to the representation studied in this section.\par
 
% END S6.p1
 
% BEGIN S6.SS1
 
\par\medskip\noindent\textbf{6.1 Automata-Based Finite Representation}\par
 
% BEGIN S6.SS1.p1
 
The framework we introduced is designed to be general, including the representation of the reward function. Here, we study one possible representation, which is based on automata over finite and infinite words. We need both finite automata to model terminating components, as well as $\omega$-automata to account for infinite behaviors of reactive systems.\par
 
% END S6.SS1.p1
 
% BEGIN S6.SS1.p2
 
We consider functions $\mathcal{F}:\Sigma^{\infty}\to\mathbb{R}$ such that $\mathit{Vals}(\mathcal{F})$ is finite, and which have a finite representation as a set of pairs of automata $\Phi=\{(\mathcal{B}_{v},\mathcal{N}_{v})\mid v\in\mathit{Vals}(\mathcal{F})\}$, where for each $v\in\mathit{Vals}(\mathcal{F})$, $\mathcal{B}_{v}$ is an NBA, and $\mathcal{N}_{v}$ is an NFA, and $\mathcal{L}(\mathcal{B}_{v})\cup\mathcal{L}(\mathcal{N}_{v})=\mathcal{F}^{=v}$. Intuitively, for each of the finitely many possible values of $\mathcal{F}$, the set of finite (resp. infinite) words that $\mathcal{F}$ maps to that value is given as a NFA (resp. NBA). We denote with $\llbracket\Phi\rrbracket$ the function $\mathcal{F}:\Sigma^{\infty}\to\mathbb{R}$ represented by $\Phi$. For a reward interface $P=(A_{P},\mathcal{F}_{P})$ where the function $\mathcal{F}_{P}$ is given in a finite representation $\Phi_{P}$, we abuse notation and write $P=(A_{P},\Phi_{P})$, and $\mathit{Vals}(\Phi_{P})$ instead of $\mathit{Vals}(\llbracket\Phi_{P}\rrbracket)$. The size of $\Phi_{P}$ is given by $|\Phi_{P}|=\sum_{v\in\mathit{Vals}(\mathcal{F}_{P})}(|\mathcal{B}_{v}|+|\mathcal{N}_{v}|)$.\par
 
% END S6.SS1.p2
 
% END S6.SS1
 
% BEGIN S6.SS2
 
\par\medskip\noindent\textbf{6.2 Decision Problems and Automata-Based Algorithms}\par
 
% BEGIN S6.SS2.p1
 
Assuming reward functions $\llbracket\Phi\rrbracket$ given as defined above, with pairs of automata for each value, we now study the main decision problems for reward interfaces.\par
 
% END S6.SS2.p1
 
% BEGIN Thmtheorem26
 
\par\medskip\noindent\textbf{Theorem 26 (Checking Compatibility of Reward Interfaces).}\par
 
% BEGIN Thmtheorem26.p1
 
For a given finitely represented reward function $\llbracket\Phi\rrbracket:(\Sigma_{P\otimes Q}^{Obs})^{\infty}\to\mathbb{R}$, checking if $P=(A_{P},\Phi_{P})$ and $Q=(A_{Q},\Phi_{Q})$ are $\llbracket\Phi\rrbracket$-compatible can be done in double exponential time.\par
 
% END Thmtheorem26.p1
 
% END Thmtheorem26
 
% BEGIN S6.SS2.SSS0.P0.SPx1
 
\par\medskip\noindent\textbf{Proof Sketch.}\par
 
% BEGIN S6.SS2.SSS0.P0.SPx1.p1
 
We reduce the second condition of Definition 15 to checking language inclusion for tree automata, by constructing an automaton capturing the set of joint implementations $S_{P}\parallel S_{Q}$ s.t. $S_{P}\in\mathsf{Imp}(P)$ and $S_{Q}\in\mathsf{Imp}(Q)$. To this end, we construct a deterministic interface automaton from $A_{P}\parallel A_{Q}$, which incurs an exponential blowup in the worst case. We construct two universal co-Büchi word automata, one for implementations $S_{P}\parallel S_{Q}$, and one for good-enough sequences w.r.t. $\llbracket\Phi\rrbracket$, that are converted to universal co-Büchi tree automata. Checking language inclusion can be done in exponential time, thus in total we can check reward interface compatibility in time double exponential in $|A_{P}|\cdot|A_{Q}|\cdot|\Phi_{P}|\cdot|\Phi_{Q}|$. $\hfill\blacktriangleleft$\par
 
% END S6.SS2.SSS0.P0.SPx1.p1
 
% BEGIN Thmtheorem27
 
\par\medskip\noindent\textbf{Theorem 27 (Checking Reward Interface Refinement).}\par
 
% BEGIN Thmtheorem27.p1
 
Checking if $(A_{Q},\Phi_{Q})\preceq(A_{P},\Phi_{P})$ can be done in exponential time.\par
 
% END Thmtheorem27.p1
 
% END Thmtheorem27
 
% END S6.SS2.SSS0.P0.SPx1
 
% BEGIN S6.SS2.SSS0.P0.SPx2
 
\par\medskip\noindent\textbf{Proof Sketch.}\par
 
% BEGIN S6.SS2.SSS0.P0.SPx2.p1
 
Checking $A_{Q}\preceq A_{P}$ is done in polynomial time [14]. To find a candidate function $r_{PQ}$ according to condition 1 of Definition 21, we iteratively match values of $\mathit{Vals}(\Phi_{P})$ to those in $\mathit{Vals}(\Phi_{Q})$. We compare sets of hopeful input sequences for the respective pairs of values in exponential time, by checking language inclusion. If a candidate $r_{PQ}$ was constructed, we verify condition 2 of Definition 21, by checking that for every complete trace achieving value $r_{PQ}(v)$ on $\llbracket\Phi_{Q}\rrbracket$, the corresponding traces achieve value $v$ for $\llbracket\Phi_{P}\rrbracket$. This is done by checking language emptiness. The overall procedure runs in exponential time. $\hfill\blacktriangleleft$\par
 
% END S6.SS2.SSS0.P0.SPx2.p1
 
% BEGIN Thmtheorem28
 
\par\medskip\noindent\textbf{Theorem 28 (Implementations of a Reward Interface).}\par
 
% BEGIN Thmtheorem28.p1
 
Checking if $S\in\mathsf{Imp}((A_{P},\Phi_{P}))$ for an interface automaton $S$ can be done in polynomial time. Checking if $\mathsf{Imp}(P)\neq\emptyset$ for $P=(A_{P},\Phi_{P})$ can be done in time exponential in $|A_{P}|\cdot|\Phi_{P}|$.\par
 
% END Thmtheorem28.p1
 
% END Thmtheorem28
 
% END S6.SS2.SSS0.P0.SPx2
 
% BEGIN S6.SS2.SSS0.P0.SPx3
 
\par\medskip\noindent\textbf{Proof Sketch.}\par
 
% BEGIN S6.SS2.SSS0.P0.SPx3.p1
 
From $\Phi_{P}$, we construct an NBA whose language contains all counterexamples a good-enough implementation must not produce, whose size is polynomial in $|\Phi_{P}|$. We check for language emptiness of its product with $S$, which can be done in polynomial time.\par
 
% END S6.SS2.SSS0.P0.SPx3.p1
 
% BEGIN S6.SS2.SSS0.P0.SPx3.p2
 
For checking the existence of an implementation, we construct from the automaton accepting counterexamples a deterministic parity automaton for the positive case, incurring in the worst case an exponential blow-up in the number of states. We intersect the result with a deterministic automaton obtained from $A_{P}$ to restrict the language to traces with corresponding executions on $A_{P}$. From the combined automaton we construct a parity game and solve it. If there exists a winning strategy, it gives us a possible implementation of $P$. If there is no such strategy, then there exists no implementation. Since the number of states in the game is exponential in $|A_{P}|\cdot|\Phi_{P}|$, the overall check can be done in exponential time. $\hfill\blacktriangleleft$\par
 
% END S6.SS2.SSS0.P0.SPx3.p2
 
% END S6.SS2.SSS0.P0.SPx3
 
% END S6.SS2
 
% BEGIN S6.SS3
 
\par\medskip\noindent\textbf{6.3 Discussion on Reward Functions and their Representation}\par
 
% BEGIN S6.SS3.p1
 
The automata-based finite representation of reward functions allows system designers to express a quantitative language with a finite set of possible values in a convenient way, by describing the languages mapped to individual values as finite automata, providing a ranking for sets of execution traces with respect to how desirable these executions are.\par
 
% END S6.SS3.p1
 
% BEGIN S6.SS3.p2
 
This finite representation enables higher-level specification languages for quantitative requirements, such as LTL[$\mathcal{F}$] [1], which is a temporal logic with quantitative features, and the derived logic LTLf[$\mathcal{F}$] [5] over finite traces. Such logic formulas can be seen as functions mapping infinite and finite words, respectively, to values in $\mathbb{R}$. The range of values for LTL[$\mathcal{F}$] specifications is finite [1]. Following the construction in [1], we can translate, in exponential time, LTL[$\mathcal{F}$] and LTLf[$\mathcal{F}$] formulas into the above automata representation.\par
 
% END S6.SS3.p2
 
% BEGIN S6.SS3.p3
 
A reward function can also capture a quantitative language $L_{\mathcal{W}}:\Sigma^{\infty}\to\mathbb{R}$, expressed with a weighted automaton $\mathcal{W}=\langle Q,q_{I},\Sigma,\delta,\gamma\rangle$, where $\gamma:\delta\to\mathbb{Q}$ is a weight function, and a value function $\mathsf{V}:\mathbb{Q}^{\infty}\to\mathbb{R}$ [8]. These automata can serve as a concise specification in the design process. To use the finite representation above and enable the respective algorithms, we restrict the instances of $\mathsf{V}$ we consider to those where $\mathit{Vals}$ finite. For the value functions $\mathit{Last}$ or $\mathit{Max}$ on finite sequences, and $\mathit{Sup}$, $\mathit{LimSup}$ or $\mathit{LimInf}$ on infinite sequences, we can construct in at most polynomial time [8], from a quantitative automaton $\mathcal{W}$, individual (Boolean) automata $\mathcal{W}^{\sim v}$ for threshold $v$ and $\sim\in\{>,\geq,=,\leq,<\}$, such that $L_{\mathcal{W}^{\sim v}}=\{\sigma\in\Sigma^{\infty}\mid L_{\mathcal{W}}(\sigma)\sim v\}$. We can thus translate a pair of quantitative automata with those value functions into a set of automata $\{(\mathcal{B}_{v},\mathcal{N}_{v})\mid v\in\mathit{Vals}\}$. While our theory supports more expressive classes of quantitative languages, such as $\mathit{LimAvg}$ or $\mathit{Disc}$, there the language inclusion problem becomes undecidable for nondeterministic automata [8] and deterministic automata lack many essential closure properties [7]. In the future, we plan to investigate decidable subclasses, in order to extend algorithmic support to more expressive representations of reward functions.\par
 
% END S6.SS3.p3
 
% END S6.SS3
 
% END S6


% BEGIN A1
 
\par\medskip\noindent\textbf{Appendix A Appendix: Compatibility of Reward Interfaces}\par
 
% BEGIN A1.p1
 
% BEGIN Thmtheorem18
 
\par\medskip\noindent\textbf{Proposition 18 (${{\mathcal{F}_{P}}\mathbin{\triangledown_{\mathit{comb}}}{\mathcal{F}_{Q}}}$-Compatibility). [Restated, see original statement.]}\par
 
% BEGIN Thmtheorem18.p1
 
Let $\mathit{comb}:\mathbb{R}_{-\infty}\times\mathbb{R}_{-\infty}\to\mathbb{R}_{-\infty}$. For all reward interfaces $P=(A_{P},\mathcal{F}_{P})$ and $Q=(A_{Q},\mathcal{F}_{Q})$ with $A_{P}$ and $A_{Q}$ compatible, it holds that $P$ and $Q$ are ${{\mathcal{F}_{P}}\mathbin{\triangledown_{\mathit{comb}}}{\mathcal{F}_{Q}}}$-compatible.\par
 
% END Thmtheorem18.p1
 
% END Thmtheorem18
 
% END A1.p1
 
% BEGIN A1.SS0.SSS0.P0.SPx1
 
\par\medskip\noindent\textbf{Proof.}\par
 
% BEGIN A1.SS0.SSS0.P0.SPx1.p1
 
Let $P=(A_{P},\mathcal{F}_{P})$ and $Q=(A_{Q},\mathcal{F}_{Q})$ be reward interfaces where $A_{P}$ and $A_{Q}$ are compatible. To prove that $P$ and $Q$ are ${{\mathcal{F}_{P}}\mathbin{\triangledown_{\mathit{comb}}}{\mathcal{F}_{Q}}}$-compatible, we need to show that for any pair of implementations $S_{P}\in\mathsf{Imp}(P)$ and $S_{Q}\in\mathsf{Imp}(Q)$ that are composable, it holds that $S_{P}\parallel S_{Q}$ is good enough w.r.t. ${{\mathcal{F}_{P}}\mathbin{\triangledown_{\mathit{comb}}}{\mathcal{F}_{Q}}}$. For the sake of contradiction, assume that there exist composable $S_{P}\in\mathsf{Imp}(P)$ and $S_{Q}\in\mathsf{Imp}(Q)$ such that $S_{P}\parallel S_{Q}$ is not good-enough w.r.t. ${{\mathcal{F}_{P}}\mathbin{\triangledown_{\mathit{comb}}}{\mathcal{F}_{Q}}}$. That means that for some $v\in\mathit{Vals}({{\mathcal{F}_{P}}\mathbin{\triangledown_{\mathit{comb}}}{\mathcal{F}_{Q}}})$ there exists an input sequence $\sigma_{E,I}\in(\Sigma_{P\otimes Q}^{E,I})^{\infty}$ that witnesses this violation, that is,\par
  
\par\noindent  $\blacksquare$ 
% BEGIN A1.I1.i1.p1
 
$\sigma_{E,I}\in\mathit{Hopeful}({{\mathcal{F}_{P}}\mathbin{\triangledown_{\mathit{comb}}}{\mathcal{F}_{Q}}},v,\Sigma_{P\otimes Q}^{E,I})$, and\par
 
% END A1.I1.i1.p1
 
 
\par\noindent  $\blacksquare$ 
% BEGIN A1.I1.i2.p1
 
for some $\sigma\in\mathsf{Traces}(S_{P}\parallel S_{Q},\sigma_{E,I})$, ${{\mathcal{F}_{P}}\mathbin{\triangledown_{\mathit{comb}}}{\mathcal{F}_{Q}}}(\sigma)$ is undef. or ${{\mathcal{F}_{P}}\mathbin{\triangledown_{\mathit{comb}}}{\mathcal{F}_{Q}}}(\sigma)<v$.\par
 
% END A1.I1.i2.p1
 
  
By the first item above, together with the definition of hopeful sequences (Definition 11) and the definition of ${{\mathcal{F}_{P}}\mathbin{\triangledown_{\mathit{comb}}}{\mathcal{F}_{Q}}}$ (Definition 17), there exist composable $S_{P}^{\prime}\in\mathsf{Imp}(P)$ and $S_{Q}^{\prime}\in\mathsf{Imp}(Q)$ and $\sigma^{\prime}\in\mathsf{Traces}(S_{P}^{\prime}\parallel S_{Q}^{\prime})$ such that $\sigma^{\prime}|_{\Sigma_{P\otimes Q}^{E,I}}=\sigma_{E,I}$ and\par
  
\par  \quad  \[\begin{array}[]{lll}\inf\{\mathit{comb}(\mathcal{F}_{P}(\sigma^{\prime\prime}|_{\Sigma_{P}^{Obs}}),\mathcal{F}_{Q}(\sigma^{\prime\prime}|_{\Sigma_{Q}^{Obs}}))&\mid&S_{P}^{\prime\prime}\in\mathsf{Imp}(P)\text{ and }S_{Q}^{\prime\prime}\in\mathsf{Imp}(Q)\text{ composable},\\
&&\sigma^{\prime\prime}\in\mathsf{Traces}(S_{P}^{\prime\prime}\parallel S_{Q}^{\prime\prime},\sigma_{E,I})\}\geq v.\end{array}\]\quad  \quad  
  
% END A1.SS0.SSS0.P0.SPx1.p1
 
% BEGIN A1.SS0.SSS0.P0.SPx1.p2
 
Since $\sigma\in\mathsf{Traces}(S_{P}\parallel S_{Q},\sigma_{E,I})$, by the definition of ${{\mathcal{F}_{P}}\mathbin{\triangledown_{\mathit{comb}}}{\mathcal{F}_{Q}}}$, we have that ${{\mathcal{F}_{P}}\mathbin{\triangledown_{\mathit{comb}}}{\mathcal{F}_{Q}}}(\sigma)$ is defined. Thus, by the second item we have that ${{\mathcal{F}_{P}}\mathbin{\triangledown_{\mathit{comb}}}{\mathcal{F}_{Q}}}(\sigma)<v$. This, together with the definition of ${{\mathcal{F}_{P}}\mathbin{\triangledown_{\mathit{comb}}}{\mathcal{F}_{Q}}}$ implies that\par
  
\par  \quad  \[\begin{array}[]{lll}\inf\{\mathit{comb}(\mathcal{F}_{P}(\sigma^{\prime\prime}|_{\Sigma_{P}^{Obs}}),\mathcal{F}_{Q}(\sigma^{\prime\prime}|_{\Sigma_{Q}^{Obs}}))&\mid&S_{P}^{\prime\prime}\in\mathsf{Imp}(P)\text{ and }S_{Q}^{\prime\prime}\in\mathsf{Imp}(Q)\text{ composable},\\
&&\sigma^{\prime\prime}\in\mathsf{Traces}(S_{P}^{\prime\prime}\parallel S_{Q}^{\prime\prime},\sigma_{E,I})\}<v.\end{array}\]\quad  \quad  
  
This is a contradiction, which concludes the proof. $\hfill\blacktriangleleft$\par
 
% END A1.SS0.SSS0.P0.SPx1.p2
 
% BEGIN A1.SS0.SSS0.P0.SPx1.p3
 
% BEGIN Thmtheorem19
 
\par\medskip\noindent\textbf{Proposition 19 (Quality of the Reward Function Composition). [Restated, see original statement.]}\par
 
% BEGIN Thmtheorem19.p1
 
Consider reward interfaces $P=(A_{P},\mathcal{F}_{P})$ and $Q=(A_{Q},\mathcal{F}_{Q})$ with $A_{P}$ and $A_{Q}$ compatible, and a reward function $\mathcal{F}$ such that $P$ and $Q$ are $\mathcal{F}$-compatible. Then, $\mathsf{Imp}(A_{P}\parallel A_{Q},{{\mathcal{F}_{P}}\mathbin{\triangledown_{\mathit{comb}}}{\mathcal{F}_{Q}}})\subseteq\mathsf{Imp}(A_{P}\parallel A_{Q},\mathcal{F})$.\par
 
% END Thmtheorem19.p1
 
% END Thmtheorem19
 
% END A1.SS0.SSS0.P0.SPx1.p3
 
% END A1.SS0.SSS0.P0.SPx1
 
% BEGIN A1.SS0.SSS0.P0.SPx2
 
\par\medskip\noindent\textbf{Proof.}\par
 
% BEGIN A1.SS0.SSS0.P0.SPx2.p1
 
Let $P=(A_{P},\mathcal{F}_{P})$ and $Q=(A_{Q},\mathcal{F}_{Q})$ be reward interfaces with automata $A_{P}$ and $A_{Q}$ that are compatible. Let $\mathcal{F}$ be a reward function such that $P$ and $Q$ are $\mathcal{F}$-compatible.\par
 
% END A1.SS0.SSS0.P0.SPx2.p1
 
% BEGIN A1.SS0.SSS0.P0.SPx2.p2
 
We have to show that for every $S\in\mathsf{Imp}(A_{P}\parallel A_{Q},{{\mathcal{F}_{P}}\mathbin{\triangledown_{\mathit{comb}}}{\mathcal{F}_{Q}}})$ it holds that $S\in\mathsf{Imp}(A_{P}\parallel A_{Q},\mathcal{F})$. For the sake of contradiction, suppose that there exists $S\in\mathsf{Imp}(A_{P}\parallel A_{Q},{{\mathcal{F}_{P}}\mathbin{\triangledown_{\mathit{comb}}}{\mathcal{F}_{Q}}})$ such that $S\not\in\mathsf{Imp}(A_{P}\parallel A_{Q},\mathcal{F})$. Since $S\preceq A_{P}\parallel A_{Q}$, this means that $S$ is not good-enough with respect to $\mathcal{F}$. Let $\sigma\in\mathsf{Traces}(S)$ be the trace witnessing this violation. Thus, $\sigma|_{\Sigma_{P\otimes Q}^{E,I}}$ is $(\mathcal{F},v)$-hopeful for some $v\in\mathbb{R}_{-\infty}$ and $\mathcal{F}(\sigma|_{\Sigma_{P\otimes Q}^{Obs}})$ is either undefined or strictly smaller than $v$. We will now show that this assumption leads to contradiction.\par
 
% END A1.SS0.SSS0.P0.SPx2.p2
 
% BEGIN A1.SS0.SSS0.P0.SPx2.p3
 
Let $\mathcal{I}=\{S_{P}\parallel S_{Q}\mid S_{P}\in\mathsf{Imp}(P),S_{Q}\in\mathsf{Imp}(Q),S_{P}\text{ and }S_{Q}\text{ composable}\}$ be the set of composed implementations of $P$ and $Q$. Since implementations do not restrict the allowed inputs, we have that there exists $S^{\prime}\in\mathcal{I}$ and $\sigma^{\prime}\in\mathsf{Traces}(S^{\prime})$ such that $\sigma^{\prime}|_{\Sigma_{P\otimes Q}^{E,I}}=\sigma|_{\Sigma_{P\otimes Q}^{E,I}}$. This implies that ${{\mathcal{F}_{P}}\mathbin{\triangledown_{\mathit{comb}}}{\mathcal{F}_{Q}}}(\sigma|_{\Sigma_{P\otimes Q}^{E,I}})$ is defined. Thus, since $S$ is good-enough with respect to ${{\mathcal{F}_{P}}\mathbin{\triangledown_{\mathit{comb}}}{\mathcal{F}_{Q}}}$, we have that ${{\mathcal{F}_{P}}\mathbin{\triangledown_{\mathit{comb}}}{\mathcal{F}_{Q}}}(\sigma|_{\Sigma_{P\otimes Q}^{Obs}})$ must be defined as well. By the definition of ${{\mathcal{F}_{P}}\mathbin{\triangledown_{\mathit{comb}}}{\mathcal{F}_{Q}}}$, this implies that there exist $S^{\prime\prime}\in\mathcal{I}$ and $\sigma^{\prime\prime}\in\mathsf{Traces}(S^{\prime\prime})$ such that $\sigma^{\prime\prime}|_{\Sigma_{P\otimes Q}^{Obs}}=\sigma|_{\Sigma_{P\otimes Q}^{Obs}}$. Now, since $\sigma^{\prime\prime}|_{\Sigma_{P\otimes Q}^{E,I}}=\sigma|_{\Sigma_{P\otimes Q}^{E,I}}$, and by Proposition 18 $S^{\prime\prime}$ is good-enough with respect to $\mathcal{F}$, it must be the case that $\mathcal{F}(\sigma^{\prime\prime}|_{\Sigma_{P\otimes Q}^{Obs}})\geq v$. Since $\sigma^{\prime\prime}|_{\Sigma_{P\otimes Q}^{Obs}}=\sigma|_{\Sigma_{P\otimes Q}^{Obs}}$, it also holds that $\mathcal{F}(\sigma|_{\Sigma_{P\otimes Q}^{Obs}})\geq v$, which is the desired contradiction. $\hfill\blacktriangleleft$\par
 
% END A1.SS0.SSS0.P0.SPx2.p3
 
% BEGIN Thmtheorem29
 
\par\medskip\noindent\textbf{Proposition 29 (Composition Associativity).}\par
 
% BEGIN Thmtheorem29.p1
 
Let $P=(A_{P},\mathcal{F}_{P}),Q=(A_{Q},\mathcal{F}_{Q})$ and $R=(A_{R},\mathcal{F}_{R})$ be reward interfaces with pairwise compatible interface automata. If the function $\mathit{comb}$ is symmetric, then ${{\mathcal{F}_{P}}\mathbin{\triangledown_{\mathit{comb}}}{\mathcal{F}_{Q}}}={{\mathcal{F}_{Q}}\mathbin{\triangledown_{\mathit{comb}}}{\mathcal{F}_{P}}}$. If $\mathit{comb}$ is associative, monotonically increasing and continuous, then the reward function composition using $\mathit{comb}$ is associative, that is, $({{{{\mathcal{F}_{P}}\mathbin{\triangledown_{\mathit{comb}}}{\mathcal{F}_{Q})}}}\mathbin{\triangledown_{\mathit{comb}}}{\mathcal{F}_{R}}}={{\mathcal{F}_{P}}\mathbin{\triangledown_{\mathit{comb}}}{{{(\mathcal{F}_{Q}}\mathbin{\triangledown_{\mathit{comb}}}{\mathcal{F}_{R}}}}})$.\par
 
% END Thmtheorem29.p1
 
% END Thmtheorem29
 
% END A1.SS0.SSS0.P0.SPx2
 
% BEGIN A1.SS0.SSS0.P0.SPx3
 
\par\medskip\noindent\textbf{Proof.}\par
 
% BEGIN A1.SS0.SSS0.P0.SPx3.p1
 
Let $P=(A_{P},\mathcal{F}_{P}),Q=(A_{Q},\mathcal{F}_{Q})$ and $R=(A_{R},\mathcal{F}_{R})$ be reward interfaces with pairwise compatible interface automata. Clearly, $A_{P}\parallel A_{Q}=A_{Q}\parallel A_{P}$. The associativity of composition for compatible interface automata is established in [14], thus $(A_{P}\parallel A_{Q})\parallel A_{R}=A_{P}\parallel(A_{Q}\parallel A_{R})$. Let $A_{PQ}:=A_{P}\parallel A_{Q}$ and $A_{QR}:=A_{Q}\parallel A_{R}$.\par
 
% END A1.SS0.SSS0.P0.SPx3.p1
 
% BEGIN A1.SS0.SSS0.P0.SPx3.p2
 
By the definition of composition, it is clearly the case that for every $\sigma\in\Sigma_{P\otimes Q}^{Obs}$ we have that ${{\mathcal{F}_{P}}\mathbin{\triangledown_{\mathit{comb}}}{\mathcal{F}_{Q}}}(\sigma)$ is defined if and only if ${{\mathcal{F}_{Q}}\mathbin{\triangledown_{\mathit{comb}}}{\mathcal{F}_{P}}}(\sigma)$ is defined. When $\mathit{comb}$ is symmetric, that, for all $a,b\in\mathbb{R}_{-\infty}$ we have that $\mathit{comb}(a,b)=\mathit{comb}(b,a)$, the definition of reward function composition (Definition 17) implies that if defined, the two values are equal.\par
 
% END A1.SS0.SSS0.P0.SPx3.p2
 
% BEGIN A1.SS0.SSS0.P0.SPx3.p3
 
Suppose that the function $\mathit{comb}$ satisfies the following conditions for all $a,b,c,d,d^{\prime}\in\mathbb{R}_{-\infty}$:\par
  
\par\noindent  $\blacksquare$ 
% BEGIN A1.I2.i1.p1
 
(Associativity) $\mathit{comb}(comb(a,b),c)=\mathit{comb}(a,comb(b,c))$.\par
 
% END A1.I2.i1.p1
 
 
\par\noindent  $\blacksquare$ 
% BEGIN A1.I2.i2.p1
 
(Monotonicity) If $d\leq d^{\prime}$, then $\mathit{comb}(d,b)\leq\mathit{comb}(d^{\prime},b)$ and $\mathit{comb}(a,d)\leq\mathit{comb}(a,d^{\prime})$.\par
 
% END A1.I2.i2.p1
 
 
\par\noindent  $\blacksquare$ 
% BEGIN A1.I2.i3.p1
 
(Continuity) For every fixed $v\in\mathbb{R}_{-\infty}$, the functions $f_{1}(x)=\mathit{comb}(x,v)$ and $f_{2}(y)=\mathit{comb}(v,y)$ are continuous on $\mathbb{R}_{-\infty}$.\par
 
% END A1.I2.i3.p1
 
  
% END A1.SS0.SSS0.P0.SPx3.p3
 
% BEGIN A1.SS0.SSS0.P0.SPx3.p4
 
Under these conditions we prove that for any sequence $\sigma\in(\Sigma_{P\otimes Q\otimes R}^{Obs})^{\infty}$ it either holds that $(({{{{\mathcal{F}_{P}}\mathbin{\triangledown_{\mathit{comb}}}{\mathcal{F}_{Q})}}}\mathbin{\triangledown_{\mathit{comb}}}{\mathcal{F}_{R}}})(\sigma)=({{\mathcal{F}_{P}}\mathbin{\triangledown_{\mathit{comb}}}{{{(\mathcal{F}_{Q}}\mathbin{\triangledown_{\mathit{comb}}}{\mathcal{F}_{R}}}}}))(\sigma)$, or both are undefined.\par
 
% END A1.SS0.SSS0.P0.SPx3.p4
 
% BEGIN A1.SS0.SSS0.P0.SPx3.p5
 
Let $\mathcal{F}_{PQ}:={{\mathcal{F}_{P}}\mathbin{\triangledown_{\mathit{comb}}}{\mathcal{F}_{Q}}}$, $\mathcal{F}_{QR}:={{\mathcal{F}_{Q}}\mathbin{\triangledown_{\mathit{comb}}}{\mathcal{F}_{R}}}$, $PQ:=(A_{PQ},\mathcal{F}_{PQ})$ and $QR:=(A_{QR},\mathcal{F}_{QR})$, and let $A:=(A_{P}\parallel A_{Q})\parallel A_{R}=A_{P}\parallel(A_{Q}\parallel A_{R})$. Let $\sigma\in(\Sigma_{P\otimes Q\otimes R}^{Obs})^{\infty}$.\par
 
% END A1.SS0.SSS0.P0.SPx3.p5
 
% BEGIN A1.SS0.SSS0.P0.SPx3.p6
 
First we show that $({{\mathcal{F}_{PQ}}\mathbin{\triangledown_{\mathit{comb}}}{\mathcal{F}_{R}}})(\sigma)$ is defined iff $({{\mathcal{F}_{P}}\mathbin{\triangledown_{\mathit{comb}}}{\mathcal{F}_{QR}}})(\sigma)$ is defined.\par
 
% END A1.SS0.SSS0.P0.SPx3.p6
 
% BEGIN A1.SS0.SSS0.P0.SPx3.p7
 
Suppose that $({{\mathcal{F}_{PQ}}\mathbin{\triangledown_{\mathit{comb}}}{\mathcal{F}_{R}}})(\sigma)$ is defined. Then, there exist composable implementations $S_{PQ}^{\prime}\in\mathsf{Imp}(PQ)$ and $S_{R}^{\prime}\in\mathsf{Imp}(R)$ and $\sigma^{\prime}\in\mathsf{Traces}(S_{PQ}^{\prime}\parallel S_{R}^{\prime})$ such that $\sigma^{\prime}|_{\Sigma_{P\otimes Q\otimes R}^{Obs}}=\sigma$. Since $S_{PQ}^{\prime}$ is good-enough with respect to $\mathcal{F}_{PQ}$, from the properties of composition of reward functions we have that there exist composable $S_{P}^{\prime}\in\mathsf{Imp}(P)$ and $S_{Q}^{\prime}\in\mathsf{Imp}(Q)$ and a trace $\sigma^{\prime\prime}\in\mathsf{Traces}(S_{P}^{\prime}\parallel S_{Q}^{\prime})$ such that $\sigma^{\prime\prime}|_{\Sigma_{P\otimes Q}^{Obs}}=\sigma^{\prime}|_{\Sigma_{P\otimes Q}^{Obs}}$. By the choice of $\sigma^{\prime}$ and $\sigma^{\prime\prime}$ we have that there exists $\sigma^{\prime\prime\prime}\in\mathsf{Traces}((S_{P}^{\prime}\parallel S_{Q}^{\prime})\parallel S_{R}^{\prime})$ such that $\sigma^{\prime\prime\prime}|_{\Sigma_{P\otimes Q\otimes R}^{Obs}}=\sigma$. Since $\mathsf{Traces}((S_{P}^{\prime}\parallel S_{Q}^{\prime})\parallel S_{R}^{\prime})=\mathsf{Traces}(S_{P}^{\prime}\parallel(S_{Q}^{\prime}\parallel S_{R}^{\prime}))$, we can use $S_{P}^{\prime}$, $S_{Q}^{\prime}\parallel S_{R}^{\prime}$ and $\sigma^{\prime\prime\prime}$ as a witness to the fact that $({{\mathcal{F}_{P}}\mathbin{\triangledown_{\mathit{comb}}}{\mathcal{F}_{QR}}})(\sigma)$ is defined. The other direction of the implication is shown analogously, concluding the first part of the proof.\par
 
% END A1.SS0.SSS0.P0.SPx3.p7
 
% BEGIN A1.SS0.SSS0.P0.SPx3.p8
 
It remains to show that when both functions are defined on $\sigma$, then their values are equal.\par
 
% END A1.SS0.SSS0.P0.SPx3.p8
 
% BEGIN A1.SS0.SSS0.P0.SPx3.p9
 
By definition, we have $({{\mathcal{F}_{PQ}}\mathbin{\triangledown_{\mathit{comb}}}{\mathcal{F}_{R}}})(\sigma)=\inf\{\mathit{comb}(\mathcal{F}_{PQ}(\sigma^{\prime}|_{\Sigma_{P\otimes Q}^{Obs}}),\mathcal{F}_{R}(\sigma^{\prime}|_{\Sigma_{R}^{Obs}}))|\\
S^{\prime}_{PQ}\in\mathsf{Imp}(PQ),S^{\prime}_{R}\in\mathsf{Imp}(R),\text{composable},\sigma^{\prime}\in\mathsf{Traces}(S^{\prime}_{PQ}\parallel S^{\prime}_{R},\sigma|_{\Sigma_{P\otimes Q\otimes R}^{E,I}})\}.$ Applying the definition to $\mathcal{F}_{PQ}$, we replace $\mathcal{F}_{PQ}(\sigma^{\prime}|_{\Sigma_{P\otimes Q}^{Obs}})$ by $\inf\{\mathit{comb}(\mathcal{F}_{P}(\sigma^{\prime\prime}|_{\Sigma_{P}^{Obs}}),\mathcal{F}_{Q}(\sigma^{\prime\prime}|_{\Sigma_{Q}^{Obs}}))|\\
S^{\prime}_{P}\in\mathsf{Imp}(P),S^{\prime}_{Q}\in\mathsf{Imp}(Q),\text{composable},\sigma^{\prime\prime}\in\mathsf{Traces}(S^{\prime}_{P}\parallel S^{\prime}_{Q},\sigma^{\prime}|_{\Sigma_{P\otimes Q}^{E,I}})\}.$\par
 
% END A1.SS0.SSS0.P0.SPx3.p9
 
% BEGIN A1.SS0.SSS0.P0.SPx3.p10
 
Due to the monotonicity and continuity properties of $\mathit{comb}$, we get $({{\mathcal{F}_{PQ}}\mathbin{\triangledown_{\mathit{comb}}}{\mathcal{F}_{R}}})(\sigma)=\inf\{\mathit{comb}(\mathit{comb}(\mathcal{F}_{P}(\sigma^{\prime\prime}|_{\Sigma_{P}^{Obs}}),\mathcal{F}_{Q}(\sigma^{\prime\prime}|_{\Sigma_{Q}^{Obs}})),\mathcal{F}_{R}(\sigma^{\prime}|_{\Sigma_{R}^{Obs}}))\mid S^{\prime}_{PQ}\in\mathsf{Imp}(PQ),S^{\prime}_{R}\in\mathsf{Imp}(R),\\
\text{composable},\sigma^{\prime}\in\mathsf{Traces}(S^{\prime}_{PQ}\parallel S^{\prime}_{R},\sigma|_{\Sigma_{P\otimes Q\otimes R}^{E,I}}),S^{\prime}_{P}\in\mathsf{Imp}(P),S^{\prime}_{Q}\in\mathsf{Imp}(Q),\text{composable},\sigma^{\prime\prime}\in\mathsf{Traces}(S^{\prime}_{P}\parallel S^{\prime}_{Q},\sigma^{\prime}|_{\Sigma_{P\otimes Q}^{E,I}})\}.$ Applying reasoning similar to that in the first part of the proof, we get $({{\mathcal{F}_{PQ}}\mathbin{\triangledown_{\mathit{comb}}}{\mathcal{F}_{R}}})(\sigma)=\inf\{\mathit{comb}(\mathit{comb}(\mathcal{F}_{P}(\sigma^{\prime\prime\prime}|_{\Sigma_{P}^{Obs}}),\mathcal{F}_{Q}(\sigma^{\prime\prime\prime}|_{\Sigma_{Q}^{Obs}})),\mathcal{F}_{R}(\sigma^{\prime\prime\prime}|_{\Sigma_{R}^{Obs}}))\mid S^{\prime}_{P}\in\mathsf{Imp}(P),S^{\prime}_{Q}\in\mathsf{Imp}(Q),\text{composable},S^{\prime}_{R}\in\mathsf{Imp}(R)\text{ composable with }S^{\prime}_{P}\parallel S^{\prime}_{Q},\sigma^{\prime\prime\prime}\in\mathsf{Traces}(S^{\prime}_{P}\parallel S^{\prime}_{Q}\parallel S^{\prime}_{R},\sigma|_{\Sigma_{P\otimes Q\otimes R}^{E,I}})\}.$ Using the associativity of $\mathit{comb}$ we replace the term $\mathit{comb}(\mathit{comb}(\mathcal{F}_{P}(\sigma^{\prime\prime\prime}|_{\Sigma_{P}^{Obs}}),\mathcal{F}_{Q}(\sigma^{\prime\prime\prime}|_{\Sigma_{Q}^{Obs}})),\mathcal{F}_{R}(\sigma^{\prime\prime\prime}|_{\Sigma_{R}^{Obs}}))$ by $\mathit{comb}(\mathcal{F}_{P}(\sigma^{\prime\prime\prime}|_{\Sigma_{P}^{Obs}}),\mathit{comb}(\mathcal{F}_{Q}(\sigma^{\prime\prime\prime}|_{\Sigma_{Q}^{Obs}}),\mathcal{F}_{R}(\sigma^{\prime\prime\prime}|_{\Sigma_{R}^{Obs}}))).$\par
 
% END A1.SS0.SSS0.P0.SPx3.p10
 
% BEGIN A1.SS0.SSS0.P0.SPx3.p11
 
Additionally, reorganizing the composition, we obtain $({{\mathcal{F}_{PQ}}\mathbin{\triangledown_{\mathit{comb}}}{\mathcal{F}_{R}}})(\sigma)=\\
\inf\{\mathit{comb}(\mathcal{F}_{P}(\sigma^{\prime\prime\prime}|_{\Sigma_{P}^{Obs}}),\mathit{comb}(\mathcal{F}_{Q}(\sigma^{\prime\prime\prime}|_{\Sigma_{Q}^{Obs}}),\mathcal{F}_{R}(\sigma^{\prime\prime\prime}|_{\Sigma_{R}^{Obs}})))\mid
S^{\prime}_{P}\in\mathsf{Imp}(P),\\
S^{\prime}_{Q}\parallel S^{\prime}_{R}\in\mathsf{Imp}(QR),\text{composable},\sigma^{\prime\prime\prime}\in\mathsf{Traces}(S^{\prime}_{P}\parallel(S^{\prime}_{Q}\parallel S^{\prime}_{R}),\sigma|_{\Sigma_{P\otimes Q\otimes R}^{E,I}})\}.$\par
 
% END A1.SS0.SSS0.P0.SPx3.p11
 
% BEGIN A1.SS0.SSS0.P0.SPx3.p12
 
Using again the properties of $\mathit{comb}$, we get $({{\mathcal{F}_{PQ}}\mathbin{\triangledown_{\mathit{comb}}}{\mathcal{F}_{R}}})(\sigma)=\\
\inf\{\mathit{comb}(\mathcal{F}_{P}(\sigma^{\prime\prime\prime}|_{\Sigma_{P}^{Obs}}),\mathcal{F}_{QR}(\sigma^{\prime\prime\prime}|_{\Sigma_{Q\otimes R}^{Obs}}))\mid S^{\prime}_{P}\in\mathsf{Imp}(P),S^{\prime}_{Q}\parallel S^{\prime}_{R}\in\mathsf{Imp}(QR),\text{composable},\\
\sigma^{\prime\prime\prime}\in\mathsf{Traces}(S^{\prime}_{P}\parallel(S^{\prime}_{Q}\parallel S^{\prime}_{R}),\sigma^{\prime}|_{\Sigma_{P\otimes Q\otimes R}^{E,I}})\}.$ This is precisely $({{\mathcal{F}_{P}}\mathbin{\triangledown_{\mathit{comb}}}{\mathcal{F}_{QR}}})(\sigma)$. $\hfill\blacktriangleleft$\par
 
% END A1.SS0.SSS0.P0.SPx3.p12
 
% END A1.SS0.SSS0.P0.SPx3
 
% END A1


% BEGIN A2
 
\par\medskip\noindent\textbf{Appendix B Appendix: Properties of Reward Interface Refinement}\par
 
% BEGIN Thmtheorem30
 
\par\medskip\noindent\textbf{Proposition 30 (Refinement as Preorder).}\par
 
% BEGIN Thmtheorem30.p1
 
For all reward interfaces $P$, $Q$ and $R$, it holds that (1) $P\preceq P$ and (2) if $Q\preceq P$, and $R\preceq Q$, then also $R\preceq P$.\par
 
% END Thmtheorem30.p1
 
% END Thmtheorem30
 
% BEGIN A2.SS0.SSS0.P0.SPx1
 
\par\medskip\noindent\textbf{Proof.}\par
 
% BEGIN A2.SS0.SSS0.P0.SPx1.p1
 
Theorem 4.1 in [14] establishes that the refinement relation between classical interface automata is a preorder. Thus, we show only that reward function refinement is a preorder.\par
 
% END A2.SS0.SSS0.P0.SPx1.p1
 
% BEGIN A2.SS0.SSS0.P0.SPx1.p2
 
Consider reward interfaces $P=(A_{P},\mathcal{F}_{P})$, $Q=(A_{Q},\mathcal{F}_{Q})$, and $R=(A_{R},\mathcal{F}_{R})$.\par
 
% END A2.SS0.SSS0.P0.SPx1.p2
 
% BEGIN A2.SS0.SSS0.P0.SPx1.p3
 
(1) is trivial, so we show only (2). Assume that $P\succeq Q$ and $Q\succeq R$.\par
 
% END A2.SS0.SSS0.P0.SPx1.p3
 
% BEGIN A2.SS0.SSS0.P0.SPx1.p4
 
Since $A_{P}\succeq A_{Q}$ and $A_{Q}\succeq A_{R}$, we have alternating simulations $\succeq_{PQ}\subseteq V_{P}\times V_{Q}$ and $\succeq_{QR}\subseteq V_{Q}\times V_{R}$ from $A_{Q}$ to $A_{P}$, and from $A_{R}$ to $A_{Q}$, respectively. Furthermore we have that $\Sigma_{P}^{E}\subseteq\Sigma_{Q}^{E}\subseteq\Sigma_{R}^{E}$, $\Sigma_{P}^{I}\subseteq\Sigma_{Q}^{I}\subseteq\Sigma_{R}^{I}$, and $\Sigma_{R}^{O}\subseteq\Sigma_{Q}^{O}\subseteq\Sigma_{P}^{O}$. For $A_{P}\succeq A_{R}$, we define the relation $\succeq_{PR}\subset V_{P}\times V_{R}$ as $\succeq_{PR}:=\{(u,v)\in V_{P}\times V_{R}\mid\exists w\in V_{Q}.\nobreak\ u\succeq_{PQ}w\land w\succeq_{QR}v\}$. The relation $\succeq_{PR}$ is an alternating simulation from $A_{R}$ to $A_{P}$. Thus, $A_{P}\succeq A_{R}$.\par
 
% END A2.SS0.SSS0.P0.SPx1.p4
 
% BEGIN A2.SS0.SSS0.P0.SPx1.p5
 
Since $\mathcal{F}_{P}\succeq\mathcal{F}_{Q}$ and $\mathcal{F}_{Q}\succeq\mathcal{F}_{R}$, there exist functions $r_{PQ}:\mathbb{R}_{-\infty}\to\mathbb{R}_{-\infty}$, and $r_{QR}:\mathbb{R}_{-\infty}\to\mathbb{R}_{-\infty}$, which satisfy the conditions of Definition 21. We define the function $r_{PR}:\mathbb{R}_{-\infty}\to\mathbb{R}_{-\infty}$ such that $r_{PR}(v)=r_{QR}(r_{PQ}(v))$ for all $v\in\mathit{Vals}(\mathcal{F}_{P})$. We now show that $r_{PR}$ satisfies the conditions of Definition 21. Let $v\in\mathit{Vals}(\mathcal{F}_{P})$.\par
  
\par\noindent  1. 
% BEGIN A2.I1.i1.p1
 
Let $\sigma_{R}^{I,E}\in(\Sigma_{R}^{E}\cup\Sigma_{R}^{I})^{\infty}$ be such that $\sigma_{R}^{I,E}|_{\Sigma_{P}}\in\mathit{Hopeful}(\mathcal{F}_{P},v,(\Sigma_{P}^{E}\cup\Sigma_{P}^{I}))$. We have to show that $\sigma_{R}^{I,E}\in\mathit{Hopeful}(\mathcal{F}_{R},r_{PR}(v),(\Sigma_{R}^{E}\cup\Sigma_{R}^{I}))$. From the properties of $r_{PQ}$, we get $\sigma_{R}^{I,E}|_{\Sigma_{Q}}\in\mathit{Hopeful}(\mathcal{F}_{Q},r_{PQ}(v),(\Sigma_{Q}^{E}\cup\Sigma_{Q}^{I}))$. From that, using the properties of $r_{QR}$, we get that $\sigma_{R}^{I,E}\ \in\mathit{Hopeful}(\mathcal{F}_{R},r_{QR}(r_{PQ}(v)),(\Sigma_{Q}^{E}\cup\Sigma_{Q}^{I}))$. Since $r_{QR}(r_{PQ}(v))=r_{PR}(v)$, this is precisely what we needed to show.\par
 
% END A2.I1.i1.p1
 
 
\par\noindent  2. 
% BEGIN A2.I1.i2.p1
 
Let $\sigma_{R}\in(\Sigma_{R}^{Obs})^{\infty}$ be such that $\mathcal{F}_{R}(\sigma_{R})\geq r_{PR}(v)$. We have to show that $\mathcal{F}_{P}(\sigma_{R}|_{\Sigma_{P}})\geq v$. Since $r_{PR}(v)=r_{QR}(r_{PQ}(v))$, from the properties of $r_{QR}$ we have that $\mathcal{F}_{Q}(\sigma_{R}|_{\Sigma_{Q}})\geq r_{PQ}(v)$. From that, by the properties of the function $r_{PQ}$ we get $\mathcal{F}_{P}((\sigma_{R}|_{\Sigma_{Q}})|_{\Sigma_{P}})\geq v$. $(\sigma_{R}|_{\Sigma_{Q}})|_{\Sigma_{P}}$ is equal to $\sigma_{R}|_{\Sigma_{P}}$ because of the subset relationships between alphabets of $P$, $Q$, and $R$, as stated above, which is precisely what we had to show.\par
 
% END A2.I1.i2.p1
 
  
$\hfill\blacktriangleleft$\par
 
% END A2.SS0.SSS0.P0.SPx1.p5
 
% BEGIN A2.SS0.SSS0.P0.SPx1.p6
 
% BEGIN Thmtheorem24
 
\par\medskip\noindent\textbf{Theorem 24 (Implementation of a Refinement). [Restated, see original statement.]}\par
 
% BEGIN Thmtheorem24.p1
 
If a reward interface $Q=(A_{Q},\mathcal{F}_{Q})$ refines a reward interface $P=(A_{P},\mathcal{F}_{P})$, then $\mathsf{Imp}(Q)\subseteq\mathsf{Imp}(P)$.\par
 
% END Thmtheorem24.p1
 
% END Thmtheorem24
 
% END A2.SS0.SSS0.P0.SPx1.p6
 
% END A2.SS0.SSS0.P0.SPx1
 
% BEGIN A2.SS0.SSS0.P0.SPx2
 
\par\medskip\noindent\textbf{Proof.}\par
 
% BEGIN A2.SS0.SSS0.P0.SPx2.p1
 
Consider reward interfaces $P=(A_{P},\mathcal{F}_{P})$ and $Q=(A_{Q},\mathcal{F}_{Q})$ such that $Q$ is a refinement of $P$. We have to show that $\mathsf{Imp}(Q)\subseteq\mathsf{Imp}(P)$. Let $S\in\mathsf{Imp}(Q)$. By Definition 13, we need to show that (a) $S\preceq A_{P}$, and (b) $S$ is good-enough w.r.t. $\mathcal{F}_{P}$.\par
 
% END A2.SS0.SSS0.P0.SPx2.p1
 
% BEGIN A2.SS0.SSS0.P0.SPx2.p2
 
Refinement of interface automata is transitive, and since $S$ is a refinement of $A_{Q}$, and $A_{Q}$ refines $A_{P}$, condition (a) is satisfied. To establish condition (b), we have to show that for every $\sigma_{E,I}\in(\Sigma_{P}^{E}\cup\Sigma_{P}^{I})^{\infty}$ with $\sigma_{E,I}\in\mathit{Hopeful}(\mathcal{F}_{P},v,\Sigma_{P}^{E}\cup\Sigma_{P}^{I})$ and every trace $\sigma\in\mathsf{Traces}_{\Sigma_{P}^{Obs}}(S,\sigma_{E,I})$ it holds that $\mathcal{F}_{P}(\sigma)$ is defined and $\mathcal{F}_{P}(\sigma)\geq v$.\par
 
% END A2.SS0.SSS0.P0.SPx2.p2
 
% BEGIN A2.SS0.SSS0.P0.SPx2.p3
 
Since $\mathcal{F}_{Q}\preceq\mathcal{F}_{P}$, there exists a function $r_{PQ}:\mathbb{R}_{-\infty}\to\mathbb{R}_{-\infty}$ that satisfies the conditions of Definition 21. Because $\sigma_{E,I}\in\mathit{Hopeful}(\mathcal{F}_{P},v,\Sigma_{P}^{E}\cup\Sigma_{P}^{I})$, and $(\Sigma_{P}^{E}\cup\Sigma_{P}^{I})\subseteq(\Sigma_{Q}^{E}\cup\Sigma_{Q}^{I})$, we know by the first property of $r_{PQ}$ that $\sigma_{E,I}\in\mathit{Hopeful}(\mathcal{F}_{Q},r_{PQ}(v),\Sigma_{Q}^{E}\cup\Sigma_{Q}^{I})$. Then, since $S\in\mathsf{Imp}(Q)$ we have that for every trace $\sigma\in\mathsf{Traces}_{\Sigma_{Q}^{Obs}}(S,\sigma_{E,I})$, $\mathcal{F}_{Q}(\sigma)$ is defined and $\mathcal{F}_{Q}(\sigma)\geq r_{PQ}(v)$. Applying Definition 21, every trace $\sigma\in\mathsf{Traces}_{\Sigma_{Q}^{Obs}}(S,\sigma_{E,I})$ is such that $\mathcal{F}_{P}(\sigma|_{\Sigma_{P}})$ is defined and $\mathcal{F}_{P}(\sigma|_{\Sigma_{P}})\geq v$. Conditions on $r_{PQ}$ guarantee that the matching sequence $\sigma|_{\Sigma_{P}}$ achieves value $v$ for $\mathcal{F}_{P}$. Since $\Sigma_{P}^{O}\supseteq\Sigma_{Q}^{O}\supseteq\Sigma_{S}^{O}$, $\{\sigma|_{\Sigma_{P}}\mid\sigma\in\mathsf{Traces}_{\Sigma_{Q}^{Obs}}(S,\sigma_{E,I})\}=\{\sigma\mid\sigma\in\mathsf{Traces}_{\Sigma_{P}^{Obs}}(S,\sigma_{E,I})\}$. Thus, it follows that for every trace $\sigma\in\mathsf{Traces}_{\Sigma_{P}^{Obs}}(S,\sigma_{E,I})$, $\mathcal{F}_{P}(\sigma)$ is defined and $\mathcal{F}_{P}(\sigma)\geq v$, which is what we had to prove. $\hfill\blacktriangleleft$\par
 
% END A2.SS0.SSS0.P0.SPx2.p3
 
% BEGIN A2.SS0.SSS0.P0.SPx2.p4
 
% BEGIN Thmtheorem25
 
\par\medskip\noindent\textbf{Theorem 25. [Restated, see original statement.]}\par
 
% BEGIN Thmtheorem25.p1
 
Consider $P=(A_{P},\mathcal{F}_{P})$ and $Q=(A_{Q},\mathcal{F}_{Q})$ with $A_{P}$ and $A_{Q}$ compatible, and let $P^{\prime}=(A_{P^{\prime}},\mathcal{F}_{P^{\prime}})$ be a non-empty refinement of $P$, with $(\Sigma_{P^{\prime}}^{I}\setminus\Sigma_{P}^{I})\cap\Sigma_{Q}^{O}=\emptyset$. Then if $\mathit{comb}$ is monotonically increasing, $(A_{P^{\prime}}\parallel A_{Q},{{\mathcal{F}_{P^{\prime}}}\mathbin{\triangledown_{\mathit{comb}}}{\mathcal{F}_{Q}}})\preceq(A_{P}\parallel A_{Q},{{\mathcal{F}_{P}}\mathbin{\triangledown_{\mathit{comb}}}{\mathcal{F}_{Q}}})$. If $P$ and $Q$ are $\mathcal{F}$-compatible for a reward function $\mathcal{F}$, then $P^{\prime}$ and $Q$ are also $\mathcal{F}$-compatible.\par
 
% END Thmtheorem25.p1
 
% END Thmtheorem25
 
% END A2.SS0.SSS0.P0.SPx2.p4
 
% BEGIN A2.SS0.SSS0.P0.SPx2.p5
 
We will differentiate Theorem 25 into Lemma 31 and Lemma 32.\par
 
% END A2.SS0.SSS0.P0.SPx2.p5
 
% BEGIN Thmtheorem31
 
\par\medskip\noindent\textbf{Lemma 31 (Substitutability).}\par
 
% BEGIN Thmtheorem31.p1
 
Let $P=(A_{P},\mathcal{F}_{P})$ be $\mathcal{F}$-compatible with $Q=(A_{Q},\mathcal{F}_{Q})$ for some function $\mathcal{F}:(\Sigma_{P\otimes Q}^{Obs})^{\infty}\to\mathbb{R}$. Then, any non-empty refinement $P^{\prime}=(A_{P^{\prime}},\mathcal{F}_{P^{\prime}})$ of $P$ with $(\Sigma_{P^{\prime}}^{I}\setminus\Sigma_{P}^{I})\cap\Sigma_{Q}^{O}=\emptyset$ is $\mathcal{F}$-compatible with $Q$.\par
 
% END Thmtheorem31.p1
 
% END Thmtheorem31
 
% END A2.SS0.SSS0.P0.SPx2
 
% BEGIN A2.SS0.SSS0.P0.SPx3
 
\par\medskip\noindent\textbf{Proof.}\par
 
% BEGIN A2.SS0.SSS0.P0.SPx3.p1
 
Since $P$ and $Q$ are $\mathcal{F}$-compatible, we know that all the following hold:\par
  
\par\noindent  $\blacksquare$ 
% BEGIN A2.I2.i1.p1
 
$A_{P}$ and $A_{Q}$ are non-empty and composable,\par
 
% END A2.I2.i1.p1
 
 
\par\noindent  $\blacksquare$ 
% BEGIN A2.I2.i2.p1
 
there exists a legal environment $A_{R}$ for $(A_{P},A_{Q})$, and\par
 
% END A2.I2.i2.p1
 
 
\par\noindent  $\blacksquare$ 
% BEGIN A2.I2.i3.p1
 
for $S_{P}\in\mathsf{Imp}(P)$ and $S_{Q}\in\mathsf{Imp}(Q)$ composable, $S_{P}\parallel S_{Q}$ is good-enough w.r.t. $\mathcal{F}$.\par
 
% END A2.I2.i3.p1
 
  
To prove $\mathcal{F}$-compatibility of $P^{\prime}$ and $Q$, we will first construct an interface automaton $A_{R^{\prime}}$ that is a legal environment for $(A_{P^{\prime}},A_{Q})$. We will then show that for any $S_{P}^{\prime}\in\mathsf{Imp}(P^{\prime})$ and any $S_{Q}\in\mathsf{Imp}(Q)$, $S_{P}^{\prime}\parallel S_{Q}$ is good-enough w.r.t. $\mathcal{F}$. By $(\Sigma_{P^{\prime}}^{I}\setminus\Sigma_{P}^{I})\cap\Sigma_{Q}^{O}=\emptyset$ and $\Sigma_{P^{\prime}}^{O}\subseteq\Sigma_{P}^{O}$, we know that $\mathsf{Shared}(A_{P^{\prime}},A_{Q})\subseteq\mathsf{Shared}(A_{P},A_{Q})$, i.e. $A_{P^{\prime}}$ does not synchronize with $A_{Q}$ on any new actions. We define $A_{R^{\prime}}=\langle V_{R},V_{R}^{init},\Sigma_{R^{\prime}}^{I},\Sigma_{R^{\prime}}^{O},\Sigma_{R}^{E},\Sigma_{R}^{H},\mathcal{T}^{\prime}_{R}\rangle$ such that:\par
  
\par\noindent  $\blacksquare$ 
% BEGIN A2.I3.i1.p1
 
$\Sigma_{R^{\prime}}^{I}=\Sigma_{R}^{I}\setminus(\Sigma_{P}^{O}\setminus\Sigma_{P^{\prime}}^{O})$, $\Sigma_{R^{\prime}}^{O}=\Sigma_{R}^{O}\cup(\Sigma_{P^{\prime}}^{I}\setminus\Sigma_{P}^{I})$,\par
 
% END A2.I3.i1.p1
 
 
\par\noindent  $\blacksquare$ 
% BEGIN A2.I3.i2.p1
 
$\mathcal{T}^{\prime}_{R}=\{(v,a,v^{\prime})|(v,a,v^{\prime})\in\mathcal{T}_{R}\land a\in\Sigma_{R^{\prime}}\}$.\par
 
% END A2.I3.i2.p1
 
  
% END A2.SS0.SSS0.P0.SPx3.p1
 
% BEGIN A2.SS0.SSS0.P0.SPx3.p2
 
That is, $A_{R^{\prime}}$ is equal to $A_{R}$ except for the sets of inputs and outputs actions towards $P^{\prime}\otimes Q$. These are modified to account for the fact that $\Sigma_{P}^{I}\subseteq\Sigma_{P^{\prime}}^{I}$ and $\Sigma_{P}^{O}\supseteq\Sigma_{P^{\prime}}^{O}$. Since none of the new outputs $(\Sigma_{P^{\prime}}^{I}\setminus\Sigma_{P}^{I})$ of $A_{R^{\prime}}$ are used in $\mathcal{T}_{R}$, functionally $A_{R^{\prime}}$ differs from $A_{R}$ only in not accepting outputs of $P$ that are removed in $P^{\prime}$, that is $(\Sigma_{P}^{O}\setminus\Sigma_{P^{\prime}}^{O})$. To prove that $A_{R^{\prime}}$ is a legal environment for $(A_{P^{\prime}},A_{Q})$ we need to show that conditions in Definition 5 are met. Meaning that $A_{R^{\prime}}$ is not empty and\par
  
\par\noindent  1. 
% BEGIN A2.I4.i1.p1
 
$\Sigma_{R^{\prime}}^{E}=\Sigma_{P^{\prime}\otimes Q}^{E}$, $\Sigma_{R^{\prime}}^{I}=\Sigma_{P^{\prime}\otimes Q}^{O}$, $\Sigma_{R^{\prime}}^{O}=\Sigma_{P^{\prime}\otimes Q}^{I}$.\par
 
% END A2.I4.i1.p1
 
 
\par\noindent  2. 
% BEGIN A2.I4.i2.p1
 
$A_{R^{\prime}}$ is composable with $A_{P^{\prime}}\otimes A_{Q}$, and $\mathsf{Illegal}(A_{P^{\prime}}\otimes A_{Q},A_{R})=\emptyset$.\par
 
% END A2.I4.i2.p1
 
 
\par\noindent  3. 
% BEGIN A2.I4.i3.p1
 
$\mathsf{Reach}((A_{P^{\prime}}\otimes A_{Q})\otimes A_{R^{\prime}})\cap(\mathsf{Illegal}(A_{P^{\prime}},A_{Q})\times V_{R^{\prime}})=\emptyset$\par
 
% END A2.I4.i3.p1
 
  
By the definition of $A_{R^{\prime}}$ and the properties of $A_{R}$, the first two conditions hold.\par
 
% END A2.SS0.SSS0.P0.SPx3.p2
 
% BEGIN A2.SS0.SSS0.P0.SPx3.p3
 
For the third condition, note that since $A_{P^{\prime}}$ is a refinement of $A_{P}$, $A_{P^{\prime}}$ will never reject an input accepted by $A_{P}$. In particular, $\mathsf{Reach}((A_{P^{\prime}}\otimes A_{Q})\otimes A_{R^{\prime}})$ will not contain illegal states from $(\mathsf{Illegal}(A_{P^{\prime}},A_{Q})\times V_{R^{\prime}})$ if $\mathsf{Reach}((A_{P}\otimes A_{Q})\otimes A_{R})$ did not, since $\mathcal{T}^{\prime}_{R}\subseteq\mathcal{T}_{R}$. Thus, the third condition also holds. Therefore, $A_{R^{\prime}}$ is a legal environment for $(P^{\prime},Q)$.\par
 
% END A2.SS0.SSS0.P0.SPx3.p3
 
% BEGIN A2.SS0.SSS0.P0.SPx3.p4
 
Now we show that for any two implementations $S_{P^{\prime}}\in\mathsf{Imp}(P^{\prime})$ and $S_{Q}\in\mathsf{Imp}(Q)$ that are composable, $S_{P^{\prime}}\parallel S_{Q}$ is good-enough w.r.t. $\mathcal{F}$.\par
 
% END A2.SS0.SSS0.P0.SPx3.p4
 
% BEGIN A2.SS0.SSS0.P0.SPx3.p5
 
As $P^{\prime}\preceq P$, by Theorem 24 we have that $S_{P^{\prime}}\in\mathsf{Imp}(P)$. Therefore, the $\mathcal{F}$-compatibility of $P$ and $Q$ implies that $S_{P^{\prime}}\parallel S_{Q}$ is good-enough w.r.t $\mathcal{F}$. $\hfill\blacktriangleleft$\par
 
% END A2.SS0.SSS0.P0.SPx3.p5
 
% BEGIN Thmtheorem32
 
\par\medskip\noindent\textbf{Lemma 32 (Monotonicity of Reward Function Composition).}\par
 
% BEGIN Thmtheorem32.p1
 
Consider reward interfaces $P=(A_{P},\mathcal{F}_{P})$ and $Q=(A_{Q},\mathcal{F}_{Q})$ with $A_{P}$ and $A_{Q}$ compatible, and let $P^{\prime}=(A_{P}^{\prime},\mathcal{F}_{P}^{\prime})$ be a non-empty refinement of $P$, with $(\Sigma_{P^{\prime}}^{I}\setminus\Sigma_{P}^{I})\cap\Sigma_{Q}^{I}=\emptyset$. If $\mathit{comb}$ is monotonically increasing, then $(A_{P}^{\prime}\parallel A_{Q},{{\mathcal{F}_{P}^{\prime}}\mathbin{\triangledown_{\mathit{comb}}}{\mathcal{F}_{Q}}})\preceq(A_{P}\parallel A_{Q},{{\mathcal{F}_{P}}\mathbin{\triangledown_{\mathit{comb}}}{\mathcal{F}_{Q}}})$.\par
 
% END Thmtheorem32.p1
 
% END Thmtheorem32
 
% END A2.SS0.SSS0.P0.SPx3
 
% BEGIN A2.SS0.SSS0.P0.SPx4
 
\par\medskip\noindent\textbf{Proof.}\par
 
% BEGIN A2.SS0.SSS0.P0.SPx4.p1
 
Let $P=(A_{P},\mathcal{F}_{P})$ and $Q=(A_{Q},\mathcal{F}_{Q})$ be such that $A_{P}$ and $A_{Q}$ are compatible, and $P^{\prime}=(A_{P}^{\prime},\mathcal{F}_{P}^{\prime})$ be non-empty, $P^{\prime}\preceq P$ and $(\Sigma_{P^{\prime}}^{I}\setminus\Sigma_{P}^{I})\cap\Sigma_{Q}^{I}=\emptyset$. We show that if $\mathit{comb}$ is monotonically increasing, then $(A_{P}^{\prime}\parallel A_{Q},{{\mathcal{F}_{P}^{\prime}}\mathbin{\triangledown_{\mathit{comb}}}{\mathcal{F}_{Q}}})\preceq(A_{P}\parallel A_{Q},{{\mathcal{F}_{P}}\mathbin{\triangledown_{\mathit{comb}}}{\mathcal{F}_{Q}}})$.\par
 
% END A2.SS0.SSS0.P0.SPx4.p1
 
% BEGIN A2.SS0.SSS0.P0.SPx4.p2
 
For interface automata, $A_{P}^{\prime}\parallel A_{Q}\preceq A_{P}\parallel A_{Q}$. We show $({{\mathcal{F}_{P}^{\prime}}\mathbin{\triangledown_{\mathit{comb}}}{\mathcal{F}_{Q}}})\!\preceq\!({{\mathcal{F}_{P}}\mathbin{\triangledown_{\mathit{comb}}}{\mathcal{F}_{Q}}})$.\par
 
% END A2.SS0.SSS0.P0.SPx4.p2
 
% BEGIN A2.SS0.SSS0.P0.SPx4.p3
 
From $\mathcal{F}_{P}^{\prime}\preceq\mathcal{F}_{P}$, there is a function $r:Vals(\mathcal{F}_{P})\to Vals(\mathcal{F}_{P}^{\prime})$ that satisfies the conditions of Definition 21. To show $({{\mathcal{F}_{P}^{\prime}}\mathbin{\triangledown_{\mathit{comb}}}{\mathcal{F}_{Q}}})\preceq({{\mathcal{F}_{P}}\mathbin{\triangledown_{\mathit{comb}}}{\mathcal{F}_{Q}}})$, we need to construct a function $r^{\prime}:\mathbb{R}_{-\infty}\to\mathbb{R}_{-\infty}$ that satisfies the conditions of Definition 21 for these functions.\par
 
% END A2.SS0.SSS0.P0.SPx4.p3
 
% BEGIN A2.SS0.SSS0.P0.SPx4.p4
 
For all $v\in Vals({{\mathcal{F}_{P}}\mathbin{\triangledown_{\mathit{comb}}}{\mathcal{F}_{Q}}})$, let $r^{\prime}(v)=\inf\{\mathit{comb}(r(\mathcal{F}_{P}(\sigma|_{\Sigma_{P}^{Obs}})),\mathcal{F}_{Q}(\sigma|_{\Sigma_{Q}^{Obs}}))\mid\\
\sigma\in(\Sigma_{P\otimes Q}^{Obs})^{\infty}\text{ and }\mathit{comb}(\mathcal{F}_{P}(\sigma|_{\Sigma_{P}^{Obs}}),\mathcal{F}_{Q}(\sigma|_{\Sigma_{Q}^{Obs}}))\geq v\}$.\par
 
% END A2.SS0.SSS0.P0.SPx4.p4
 
% BEGIN A2.SS0.SSS0.P0.SPx4.p5
 
For every value $v\in Vals({{\mathcal{F}_{P}}\mathbin{\triangledown_{\mathit{comb}}}{\mathcal{F}_{Q}}})$ we have:\par
  
\par\noindent  $\blacksquare$ 
% BEGIN A2.I5.i1.p1
 
Condition 1: Let $\sigma_{E,I}\in(\Sigma_{P^{\prime}\otimes Q}^{E,I})^{\infty}$ be such that $\sigma_{E,I}|_{\Sigma_{P\otimes Q}^{E,I}}\in\mathit{Hopeful}({{\mathcal{F}_{P}}\mathbin{\triangledown_{\mathit{comb}}}{\mathcal{F}_{Q}}},v,\Sigma_{P\otimes Q}^{E,I})$. This means that $\inf\{\mathit{comb}(\mathcal{F}_{P}(\sigma|_{\Sigma_{P}^{Obs}}),\mathcal{F}_{Q}(\sigma|_{\Sigma_{Q}^{Obs}}))\mid S_{P}\in\mathsf{Imp}(P),S_{Q}\in\mathsf{Imp}(Q)\text{ composable},\\
\sigma\in\mathsf{Traces}(S_{P}\parallel S_{Q},\sigma_{E,I})\}\geq v$.\par
 
% END A2.I5.i1.p1
 
% BEGIN A2.I5.i1.p2
 
We know that $\mathsf{Imp}(P^{\prime})\subseteq\mathsf{Imp}(P)$, so in particular $\mathsf{Imp}(A_{P}^{\prime}\parallel A_{Q},{{\mathcal{F}_{P}^{\prime}}\mathbin{\triangledown_{\mathit{comb}}}{\mathcal{F}_{Q}}})\subseteq\mathsf{Imp}(A_{P}\parallel A_{Q},{{\mathcal{F}_{P}}\mathbin{\triangledown_{\mathit{comb}}}{\mathcal{F}_{Q}}})$. That means that for ${{\mathcal{F}_{P}^{\prime}}\mathbin{\triangledown_{\mathit{comb}}}{\mathcal{F}_{Q}}}$ we will consider a subset of the implementations, and therefore not more traces when applying $\inf$, as for ${{\mathcal{F}_{P}}\mathbin{\triangledown_{\mathit{comb}}}{\mathcal{F}_{Q}}}$. This, together with the above inequality implies ${{\mathcal{F}_{P}^{\prime}}\mathbin{\triangledown_{\mathit{comb}}}{\mathcal{F}_{Q}}}(\sigma_{E,I})\geq r^{\prime}(v)$. Then it holds that any $\sigma\in\mathsf{Traces}(A_{P}\parallel A_{Q},\sigma_{E,I})$, for which ${{\mathcal{F}_{P}^{\prime}}\mathbin{\triangledown_{\mathit{comb}}}{\mathcal{F}_{Q}}}(\sigma)$ is defined, is a witness for $\sigma_{E,I}\in\mathit{Hopeful}({{\mathcal{F}_{P}^{\prime}}\mathbin{\triangledown_{\mathit{comb}}}{\mathcal{F}_{Q}}},r^{\prime}(v),\Sigma_{P^{\prime}\otimes Q}^{E,I})$.\par
 
% END A2.I5.i1.p2
 
 
\par\noindent  $\blacksquare$ 
% BEGIN A2.I5.i2.p1
 
Condition 2: Let $\sigma^{\prime}\in(\Sigma_{P^{\prime}\otimes Q}^{Obs})^{\infty}$ be such that $({{\mathcal{F}_{P}^{\prime}}\mathbin{\triangledown_{\mathit{comb}}}{\mathcal{F}_{Q}}})(\sigma^{\prime})\geq r^{\prime}(v)$. This means that $\inf\{\mathit{comb}(\mathcal{F}_{P}^{\prime}(\sigma|_{\Sigma_{P^{\prime}}^{Obs}}),\mathcal{F}_{Q}(\sigma|_{\Sigma_{Q}^{Obs}}))\mid\\
S_{P}^{\prime}\in\mathsf{Imp}(P^{\prime}),S_{Q}\in\mathsf{Imp}(Q),\text{ composable},\sigma\in\mathsf{Traces}(S_{P}^{\prime}\parallel S_{Q},\sigma^{\prime}|_{\Sigma_{P^{\prime}\otimes Q}^{E,I}})\}\geq r^{\prime}(v)$.\par
 
% END A2.I5.i2.p1
 
% BEGIN A2.I5.i2.p2
 
We have to show that $\inf\{\mathit{comb}(\mathcal{F}_{P}(\sigma|_{\Sigma_{P}^{Obs}}),\mathcal{F}_{Q}(\sigma|_{\Sigma_{Q}^{Obs}}))\mid\\
S_{P}\in\mathsf{Imp}(P),S_{Q}\in\mathsf{Imp}(Q),\text{ composable},\sigma\in\mathsf{Traces}(S_{P}\parallel S_{Q},\sigma^{\prime}|_{\Sigma_{P\otimes Q}^{E,I}})\}\geq v$.\par
 
% END A2.I5.i2.p2
 
% BEGIN A2.I5.i2.p3
 
Since $\mathcal{F}_{P}^{\prime}\preceq\mathcal{F}_{P}$ we have for the function $r$ that $\mathcal{F}_{P}^{\prime}(\sigma)\geq r(\mathcal{F}_{P}(\sigma|_{\Sigma_{P}^{Obs}}))$. Since this function $r$ is used in $r^{\prime}$, then $\mathit{comb}(r(\mathcal{F}_{P}(\sigma|_{\Sigma_{P}^{Obs}})),\mathcal{F}_{Q}(\sigma|_{\Sigma_{Q}^{Obs}}))\geq r^{\prime}(v)$, meaning that $\mathit{comb}(\mathcal{F}_{P}(\sigma|_{\Sigma_{P}^{Obs}}),\mathcal{F}_{Q}(\sigma|_{\Sigma_{Q}^{Obs}}))\geq v$ by definition of $r^{\prime}$ and monotonicity of $\mathit{comb}$. Thus, $\inf\{\mathit{comb}(\mathcal{F}_{P}(\sigma|_{\Sigma_{P}^{Obs}}),\mathcal{F}_{Q}(\sigma|_{\Sigma_{Q}^{Obs}}))\mid\\
S_{P}\in\mathsf{Imp}(P),S_{Q}\in\mathsf{Imp}(Q),\text{ composable},\sigma\in\mathsf{Traces}(S_{P}\parallel S_{Q},\sigma^{\prime}|_{\Sigma_{P\otimes Q}^{E,I}})\}\geq v$.\par
 
% END A2.I5.i2.p3
 
  
% END A2.SS0.SSS0.P0.SPx4.p5
 
% BEGIN A2.SS0.SSS0.P0.SPx4.p6
 
Thus, both conditions in Definition 21 hold for $r^{\prime}$. Hence, $({{\mathcal{F}_{P}^{\prime}}\mathbin{\triangledown_{\mathit{comb}}}{\mathcal{F}_{Q}}})\!\preceq\!({{\mathcal{F}_{P}}\mathbin{\triangledown_{\mathit{comb}}}{\mathcal{F}_{Q}}})$. $\hfill\blacktriangleleft$\par
 
% END A2.SS0.SSS0.P0.SPx4.p6
 
% BEGIN A2.SS0.SSS0.P0.SPx4.p7
 
\par
 
% END A2.SS0.SSS0.P0.SPx4.p7
 
% END A2.SS0.SSS0.P0.SPx4
 
% END A2
